\section{هفته‌ی 1: مدل شناختی چیست؟}

\begin{quote}
\textbf{فصل 1: مبانی مدل‌سازی شناختی} · درس 1 ساعت · \lr{lab} 2 ساعت
\textbf{مدل اصلی این هفته:} یک مدل ساده‌ی \lr{random walk} (گام تصادفی)
برای تصمیم‌های دوگزینه‌ای.
\end{quote}

\begin{center}\rule{0.5\linewidth}{0.5pt}\end{center}

\subsection{واژه‌های کلیدی}

واژه‌های مهم این درس با معادل فارسی‌شان (حتی اگر در درس‌های قبل آمده
باشند).

{\def\LTcaptype{none} % do not increment counter
\begin{longtable}[]{@{}
  >{\raggedright\arraybackslash}p{(\linewidth - 4\tabcolsep) * \real{0.3333}}
  >{\raggedright\arraybackslash}p{(\linewidth - 4\tabcolsep) * \real{0.3333}}
  >{\raggedright\arraybackslash}p{(\linewidth - 4\tabcolsep) * \real{0.3333}}@{}}
\toprule\noalign{}
\begin{minipage}[b]{\linewidth}\raggedright
\lr{English}
\end{minipage} & \begin{minipage}[b]{\linewidth}\raggedright
فارسی
\end{minipage} & \begin{minipage}[b]{\linewidth}\raggedright
توضیح کوتاه
\end{minipage} \\
\midrule\noalign{}
\endhead
\bottomrule\noalign{}
\endlastfoot
\lr{cognitive model} & مدل شناختی & توصیف دقیق و ریاضی یا محاسباتی از یک
فرایند ذهنی \\
\lr{random walk} & گام تصادفی & جمع شدن گام‌به‌گام شواهد تا رسیدن به یک
مرز \\
\lr{verbal theory} & نظریه‌ی کلامی & نظریه‌ای که فقط با کلمات بیان شده و
پیش‌بینی دقیق نمی‌دهد \\
\lr{Marr's levels} & سطوح تحلیل \lr{Marr} & سطح محاسباتی، الگوریتمی و
پیاده‌سازی \\
\lr{parameter} & پارامتر & کمیت قابل تنظیم مدل که بین افراد یا شرایط فرق
می‌کند \\
\lr{noise} & نویز & تغییرات تصادفی در شواهد یا در داده \\
\lr{simulation} & شبیه‌سازی & تولید داده‌ی مصنوعی با اجرای مدل \\
\lr{drift rate} & نرخ رانش & سرعت میانگین جمع شدن شواهد؛ نشانه‌ی کیفیت
شواهد \\
\lr{boundary separation} & فاصله‌ی مرزها & مقدار شواهد لازم برای پاسخ؛
نشانه‌ی احتیاط \\
\lr{evidence} & شواهد & اطلاعاتی که برای تصمیم جمع می‌شود \\
\lr{accuracy} & دقت & نسبت پاسخ‌های درست \\
\lr{response time (RT)} & زمان واکنش & فاصله‌ی ارائه‌ی محرک تا پاسخ \\
\lr{formal model} / \lr{formalize} & مدل فرمال / فرمال کردن & نوشتن
نظریه به شکل معادله یا الگوریتم \\
\lr{fit} / \lr{fitting} & برازش & تنظیم پارامترها تا مدل با داده جور
شود \\
\lr{forward} / \lr{inverse problem} & مسئله‌ی مستقیم / معکوس & شبیه‌سازی
(پارامتر \ensuremath{\leftarrow} داده) در برابر برآورد (داده \ensuremath{\leftarrow} پارامتر) \\
\lr{non-decision time} & زمان غیرتصمیمی & زمان دیدن محرک و اجرای حرکت
پاسخ \\
\lr{speed\lr{--}accuracy trade-off} & بده‌بستان سرعت و دقت & سریع‌تر پاسخ دادن
به قیمت خطای بیشتر \\
\lr{trial} & کوشش \lr{(trial)} & یک بار انجام تکلیف توسط شرکت‌کننده \\
\lr{probabilistic programming} & برنامه‌نویسی احتمالاتی & تعریف مدل با کد
و نمونه‌گیری خودکار \lr{(PyMC)} \\
\lr{temperature} & دما \lr{(temperature)} & پارامتر نویز در انتخاب
\lr{token}؛ \ensuremath{\beta} = \lr{1/T} \\
\lr{random seed} & \lr{seed} تصادفی & عدد شروع مولد اعداد تصادفی برای
تکرارپذیری \\
\lr{quantile} & چندک \lr{(quantile)} & مقداری که درصد مشخصی از داده زیر
آن است \\
\lr{large language model (LLM)} & مدل بزرگ زبانی & شبکه‌ی عصبی‌ای که کلمه‌ی
بعدی متن را پیش‌بینی می‌کند \\
\lr{token} & \lr{token} (واحد متن) & کلمه یا بخشی از کلمه که مدل زبانی
پیش‌بینی می‌کند \\
\end{longtable}
}

\begin{center}\rule{0.5\linewidth}{0.5pt}\end{center}

\subsection{اهداف یادگیری}

در پایان این هفته، دانشجویان می‌توانند:

\begin{enumerate}
\def\labelenumi{\arabic{enumi}.}
\tightlist
\item
  مدل شناختی را \textbf{تعریف کنند} و آن را از یک نظریه‌ی کلامی و از یک
  مدل صرفاً آماری داده‌ها \textbf{متمایز کنند}.
\item
  مدل‌ها را بر اساس هدف (توصیفی، فرایندی، هنجاری) و خانواده (نمادین،
  پیوندگرا یا \lr{connectionist}، تصادفی/پویا، بیزی) \textbf{دسته‌بندی
  کنند} و جایشان را در سه سطح \lr{Marr} \textbf{مشخص کنند}.
\item
  \textbf{توضیح دهند} که چطور پارامترهای آزاد، یک هسته‌ی قطعی و یک مدل
  نویز \lr{(noise model)} با هم یک مدل را به یک ماشین \emph{تولید داده}
  تبدیل می‌کنند.
\item
  یک مدل تصمیم‌گیری \lr{random walk} را در \lr{Python} \textbf{پیاده‌سازی}
  و \textbf{شبیه‌سازی کنند} و شبیه‌سازی را با پیش‌بینی‌های تحلیلی
  \textbf{بررسی کنند}.
\item
  \textbf{تفسیر کنند} که تغییر پارامترهای مدل \lr{(drift}، مرز) با کدام
  سازه‌های روان‌شناختی (کیفیت شواهد، احتیاط) متناظر است.
\end{enumerate}

\begin{center}\rule{0.5\linewidth}{0.5pt}\end{center}

\subsection{مقدمه}

یک نظریه‌ی کلامی آشنا درباره‌ی تصمیم‌های ساده را در نظر بگیرید:

\begin{quote}
\emph{«آدم‌ها وقتی تصمیم می‌گیرند، آن‌قدر شواهد جمع می‌کنند تا کافی شود و
بعد پاسخ می‌دهند.»}
\end{quote}

کامل به نظر می‌رسد، ولی سعی کنید از آن یک پیش‌بینی بیرون بکشید:

\begin{itemize}
\tightlist
\item
  اگر به شرکت‌کننده‌ها بگوییم \emph{دقیق باشید} به‌جای \emph{سریع باشید}،
  چقدر کندتر می‌شوند؟ دو برابر؟ 50 میلی‌ثانیه؟
\item
  آیا \textbf{خطاها} سریع‌تر از پاسخ‌های درست‌اند، کندتر، یا همان‌قدر؟
\item
  توزیع زمان واکنش \lr{(RT)} باید متقارن باشد یا چوله \lr{(skewed)}؟
\item
  اگر دیدن محرک سخت‌تر شود، آیا زمان واکنش همان‌طور تغییر می‌کند که با
  دستور «دقیق باشید»؟
\end{itemize}

این نظریه‌ی کلامی \textbf{به هیچ‌کدام از این سؤال‌ها نمی‌تواند جواب دهد}. هر
جواب به تصمیم‌هایی بستگی دارد که نظریه هرگز نگرفته است: «شواهد» یعنی چه؟
پیوسته می‌رسد یا تکه‌تکه؟ «کافی» یعنی چقدر؟ آستانه برای هر دو پاسخ یکی
است؟

به‌محض اینکه نظریه را به شکل یک مدل بنویسیم، مثلاً یک \textbf{\lr{random
walk}} که گام‌به‌گام بالا یا پایین می‌رود تا به یک مرز برسد، هر کدام از این
سؤال‌ها جواب مشخصی پیدا می‌کند. بعضی از این جواب‌ها \emph{غلط} از آب
درمی‌آیند، و این خوب است: مدلی که می‌تواند غلط باشد، مدلی است که می‌شود از
آن یاد گرفت. این پیام اصلی این درس است:

\begin{quote}
مدل‌های فرمال \lr{(formal)} ما را مجبور می‌کنند فرض‌هایمان را صریح کنیم،
پیش‌بینی‌های دقیق می‌دهند، و سؤال «آیا این نظریه درست است؟» را به سؤال «این
مدل داده‌ها را چقدر خوب، و چرا، توضیح می‌دهد؟» تبدیل می‌کنند.
\end{quote}

\paragraph{این موضوع کجا به کار می‌رود؟}

مدل‌هایی مثل \lr{random walk} این هفته فقط مثال کلاسی نیستند. هر روز در
پژوهش و کار عملی از آن‌ها استفاده می‌شود:

\begin{itemize}
\tightlist
\item
  \textbf{روان‌شناسی بالینی و روان‌پزشکی.} بیماران افسرده، مبتلا به
  \lr{ADHD} یا اضطراب اغلب در تکلیف‌های ساده «کندتر» هستند. مدل‌های تصمیم
  نشان می‌دهند \emph{چرا}: کیفیت پایین‌تر شواهد، احتیاط بیشتر، یا پاسخ
  حرکتی کندتر. این سه توضیح به مداخله‌های متفاوتی نیاز دارند.
\item
  \textbf{پژوهش درباره‌ی سالمندی.} افراد مسن‌تر کندتر پاسخ می‌دهند. برازش
  مدل‌های تصمیم نشان داده که بخش زیادی از این کندی از احتیاط بیشتر و کندی
  حرکت می‌آید، نه از ضعف در پردازش اطلاعات.
\item
  \textbf{علوم اعصاب.} نرخ شلیک نورون‌های بعضی نواحی مغز پیش از تصمیم
  کم‌کم بالا می‌رود تا به یک آستانه برسد؛ خیلی شبیه \lr{random walk} که به
  مرز می‌رسد. از پارامترهای مدل برای وصل کردن رفتار و فعالیت مغز استفاده
  می‌شود.
\item
  \textbf{کاربردهای عملی.} تصمیم شاهد عینی، خواندن تصویر پزشکی و کنترل
  کیفیت، همه تصمیم‌گیری با شواهد پرنویز هستند. همین مدل «شواهد چقدر خوب
  است» را از «تصمیم‌گیرنده چقدر محتاط است» جدا می‌کند.
\item
  \textbf{هوش مصنوعی.} ایده‌های مدل‌های شناختی (جمع شدن شواهد، یادگیری از
  خطا، استنباط احتمالاتی) در سیستم‌های هوش مصنوعی به کار رفته‌اند، و امروز
  مدل‌های هوش مصنوعی هم با روش‌های همین درس آزموده می‌شوند (بخش هوش مصنوعی
  و \lr{LLM}ها را ببینید).
\end{itemize}

\begin{center}\rule{0.5\linewidth}{0.5pt}\end{center}

\subsection{جزوه‌ی درس}

\paragraph{به زبان ساده}

\emph{نظریه} یک حرف را با کلمات می‌زند («آدم‌ها شواهد جمع می‌کنند تا به
اندازه‌ی کافی برسد»). \emph{مدل} همان حرف را آن‌قدر دقیق می‌زند که کامپیوتر
بتواند اجرایش کند. وقتی مدل اجرا شود، پیش‌بینی دقیق می‌دهد: آدم‌ها چند درصد
مواقع درست جواب می‌دهند، چقدر طول می‌کشد، و این دو با تغییر تکلیف چطور عوض
می‌شوند. مدل این هفته ساده‌ترین مدل تصمیم ممکن است. یک مهره روی یک خط عددی
را تصور کنید. در هر لحظه یک قدم بالا می‌رود (کمی شاهد به نفع جواب درست)
یا یک قدم پایین (کمی شاهد به نفع جواب غلط). وقتی به خط بالا برسد، فرد
درست جواب می‌دهد؛ وقتی به خط پایین برسد، خطا می‌کند. فقط سه عدد همه‌چیز را
کنترل می‌کنند: قدم‌ها چقدر به سمت بالا تمایل دارند (\textbf{\lr{drift}}،
کیفیت شواهد)، دو خط چقدر از هم فاصله دارند (\textbf{فاصله‌ی مرزها}، میزان
احتیاط فرد)، و مهره از کجا شروع می‌کند (\textbf{نقطه‌ی شروع}، سوگیری). با
چند بار اجرای این ماشین کوچک می‌بینیم که بده‌بستان سرعت و دقت خودبه‌خود
ظاهر می‌شود، بدون اینکه آن را در مدل گذاشته باشیم.

\paragraph{واژه‌های روان‌شناسی به زبان ساده}

\begin{itemize}
\tightlist
\item
  \textbf{شرکت‌کننده \lr{(participant} یا \lr{subject):}} کسی که در
  آزمایش شرکت می‌کند.
\item
  \textbf{محرک \lr{(stimulus}، جمع \lr{stimuli):}} هر چیزی که به
  شرکت‌کننده نشان داده یا پخش می‌شود؛ مثلاً یک تصویر، یک کلمه یا یک صدا.
\item
  \textbf{تکلیف \lr{(task):}} کاری که شرکت‌کننده باید انجام دهد؛ مثلاً
  «اگر نقطه‌ها به چپ می‌روند کلید چپ و اگر به راست می‌روند کلید راست را
  بزن». \textbf{تکلیف دوگزینه‌ای} دقیقاً دو پاسخ ممکن دارد.
\item
  \textbf{\lr{trial} (کوشش):} یک بار انجام تکلیف: یک محرک، یک پاسخ. هر
  آزمایش معمولاً برای هر شرکت‌کننده صدها \lr{trial} دارد.
\item
  \textbf{زمان واکنش \lr{(RT):}} فاصله‌ی زمانی از ظاهر شدن محرک تا زدن
  کلید، معمولاً بر حسب میلی‌ثانیه \lr{(ms).}
\item
  \textbf{دقت \lr{(accuracy):}} نسبت پاسخ‌های درست.
\item
  \textbf{شرط یا دست‌کاری \lr{(condition} / \lr{manipulation):}} تغییری
  که آزمایشگر عمداً ایجاد می‌کند (مثلاً دستورالعمل «سریع باش!» در برابر
  «دقیق باش!»، یا محرک آسان در برابر سخت) تا ببیند رفتار چطور تغییر
  می‌کند.
\item
  \textbf{سازه \lr{(construct):}} یک کمیت روان‌شناختی که مستقیم دیده
  نمی‌شود؛ مثلاً \emph{احتیاط}، \emph{توجه} یا \emph{قدرت حافظه}. یک
  پارامتر خوب مدل، یک سازه را \emph{اندازه می‌گیرد}.
\item
  \textbf{داده‌ی رفتاری:} چیزی که می‌توانیم ثبت کنیم: انتخاب‌ها و زمان‌های
  واکنش (نه سیگنال‌های مغزی).
\end{itemize}

\textbf{زمان‌بندی پیشنهادی (60 دقیقه):} انگیزه 10 · چیستی و چرایی 10 ·
انواع مدل + \lr{Marr} 15 · مدل به‌عنوان ماشین تولید داده + مسیر مدل‌سازی
10 · اولین مدل \lr{(random walk)} 12 · نمادگذاری و نکات اجرایی درس 3.

\subsubsection{1. مدل شناختی چیست؟}

\textbf{مدل شناختی} توصیفی صریح، معمولاً ریاضی یا محاسباتی، از فرایندها
یا بازنمایی‌هایی است که فرض می‌کنیم رفتار را تولید می‌کنند. سه جزء دارد:

{\def\LTcaptype{none} % do not increment counter
\begin{longtable}[]{@{}
  >{\raggedright\arraybackslash}p{(\linewidth - 4\tabcolsep) * \real{0.3333}}
  >{\raggedright\arraybackslash}p{(\linewidth - 4\tabcolsep) * \real{0.3333}}
  >{\raggedright\arraybackslash}p{(\linewidth - 4\tabcolsep) * \real{0.3333}}@{}}
\toprule\noalign{}
\begin{minipage}[b]{\linewidth}\raggedright
جزء
\end{minipage} & \begin{minipage}[b]{\linewidth}\raggedright
به چه سؤالی جواب می‌دهد
\end{minipage} & \begin{minipage}[b]{\linewidth}\raggedright
مثال در \lr{random walk}
\end{minipage} \\
\midrule\noalign{}
\endhead
\bottomrule\noalign{}
\endlastfoot
\textbf{فرض‌ها / ساختار} & چه سازوکاری را فرض می‌کنیم؟ & شواهد در گام‌های
گسسته بین دو مرز انباشته می‌شود \\
\textbf{پارامترها} & کدام کمیت‌ها می‌توانند بین افراد یا شرایط تغییر کنند؟
& \lr{drift} \(v\)، فاصله‌ی مرزها \(a\)، نقطه‌ی شروع \(z\) \\
\textbf{نگاشت به داده} & سازوکار چگونه پاسخ‌های قابل مشاهده تولید می‌کند؟
& مرزی که به آن می‌رسیم انتخاب را تعیین می‌کند؛ تعداد گام‌ها زمان واکنش
را \\
\end{longtable}
}

سه اصطلاح مرتبط که به‌راحتی با هم قاطی می‌شوند:

\begin{itemize}
\tightlist
\item
  \textbf{نظریه} مجموعه‌ای از ادعاهای کلامی (اغلب غیررسمی) درباره‌ی ذهن
  است.
\item
  \textbf{مدل} یک نمونه‌ی فرمال مشخص از یک نظریه است. یک نظریه را می‌شود
  به شکل‌های مختلفی فرمال کرد.
\item
  \textbf{مدل آماری} (مثل \lr{t-test} یا رگرسیون) هم مدل است، ولی مدلِ
  \emph{داده‌ها}. پارامترهایش (میانگین‌ها، شیب‌ها) معمولاً داده‌ها را توصیف
  می‌کنند، نه یک فرایند شناختی را. در مدل شناختی، پارامترها قرار است
  \textbf{کمیت‌های روان‌شناختی پنهان را اندازه بگیرند}، مثل زوال حافظه،
  احتیاط یا توجه.
\end{itemize}

جمله‌ی معروف \lr{George Box} که «همه‌ی مدل‌ها غلط‌اند، ولی بعضی‌شان مفیدند»
نگرش این درس را خلاصه می‌کند. ما هیچ‌وقت نمی‌پرسیم یک مدل \emph{درست} است
یا نه. می‌پرسیم \emph{مفید} است یا نه: آیا الگوهای مهم را نشان می‌دهد، به
داده‌ی جدید تعمیم پیدا می‌کند، و پارامترهای قابل تفسیر دارد؟

\subsubsection{2. چرا اصلاً مدل فرمال بسازیم؟}

\begin{enumerate}
\def\labelenumi{\arabic{enumi}.}
\tightlist
\item
  \textbf{دقت.} نوشتن مدل فرض‌های پنهان را آشکار می‌کند (مثل مثال بخش
  انگیزه).
\item
  \textbf{پیش‌بینی کمّی.} مدل‌ها می‌گویند \emph{چقدر}، نه فقط \emph{آیا}.
\item
  \textbf{کشف پیامدهای غیرمنتظره.} شهود ما درباره‌ی سیستم‌های پیچیده ضعیف
  است. مثلاً میانگین‌گیری از منحنی‌های یادگیری نمایی افراد، منحنی‌ای می‌سازد
  که شبیه قانون \emph{توانی} است (هفته‌ی 5).
\item
  \textbf{اندازه‌گیری.} پارامترهای برازش‌شده اندازه‌گیری سازه‌های پنهان‌اند.
  \(d'\) در نظریه‌ی تشخیص علامت (هفته‌ی 2) حساسیت را جدا از سوگیری پاسخ
  اندازه می‌گیرد؛ کاری که درصد پاسخ درست به‌تنهایی نمی‌تواند بکند.
\item
  \textbf{مقایسه‌ی نظریه‌ها.} دو نظریه‌ی کلامی که «هر دو داده‌ها را توضیح
  می‌دهند» بعد از فرمال شدن قابل رقابت با هم‌اند (هفته‌های 6\lr{--}7 و 10).
\item
  \textbf{پلی به هوش مصنوعی.} همین صورت‌بندی‌ها (مدل‌های احتمالاتی، قواعد
  یادگیری، شبکه‌های عصبی) برای ساختن عامل‌های مصنوعی هم به کار می‌روند
  (هفته‌ی 15).
\end{enumerate}

\subsubsection{3. انواع مدل‌های شناختی}

\textbf{بر اساس هدف:}

{\def\LTcaptype{none} % do not increment counter
\begin{longtable}[]{@{}
  >{\raggedright\arraybackslash}p{(\linewidth - 6\tabcolsep) * \real{0.2500}}
  >{\raggedright\arraybackslash}p{(\linewidth - 6\tabcolsep) * \real{0.2500}}
  >{\raggedright\arraybackslash}p{(\linewidth - 6\tabcolsep) * \real{0.2500}}
  >{\raggedright\arraybackslash}p{(\linewidth - 6\tabcolsep) * \real{0.2500}}@{}}
\toprule\noalign{}
\begin{minipage}[b]{\linewidth}\raggedright
نوع
\end{minipage} & \begin{minipage}[b]{\linewidth}\raggedright
سؤال
\end{minipage} & \begin{minipage}[b]{\linewidth}\raggedright
مثال
\end{minipage} & \begin{minipage}[b]{\linewidth}\raggedright
هفته‌ها
\end{minipage} \\
\midrule\noalign{}
\endhead
\bottomrule\noalign{}
\endlastfoot
\textbf{توصیفی \lr{(descriptive)}} & شکل داده‌ها \emph{چیست}؟ & قانون
توانی تمرین؛ منحنی فراموشی نمایی & 2\lr{--}7 \\
\textbf{فرایندی \lr{(process} / \lr{mechanistic)}} & ذهن \emph{چگونه}
رفتار را تولید می‌کند؟ & \lr{random walk} / \lr{drift-diffusion}؛ مدل
\lr{exemplar (GCM)} & 1، 14 \\
\textbf{هنجاری \lr{(normative} / \lr{rational)}} & یک عامل ایده‌آل
\emph{باید} چه کند؟ & مشاهده‌گر ایده‌آل بیزی برای ترکیب سرنخ‌ها & 15 \\
\end{longtable}
}

مرزها مبهم‌اند: مدل \lr{SDT} تا حدی توصیفی و تا حدی فرایندی است.

\textbf{بر اساس خانواده:}

{\def\LTcaptype{none} % do not increment counter
\begin{longtable}[]{@{}
  >{\raggedright\arraybackslash}p{(\linewidth - 4\tabcolsep) * \real{0.3333}}
  >{\raggedright\arraybackslash}p{(\linewidth - 4\tabcolsep) * \real{0.3333}}
  >{\raggedright\arraybackslash}p{(\linewidth - 4\tabcolsep) * \real{0.3333}}@{}}
\toprule\noalign{}
\begin{minipage}[b]{\linewidth}\raggedright
خانواده
\end{minipage} & \begin{minipage}[b]{\linewidth}\raggedright
ایده‌ی اصلی
\end{minipage} & \begin{minipage}[b]{\linewidth}\raggedright
مثال
\end{minipage} \\
\midrule\noalign{}
\endhead
\bottomrule\noalign{}
\endlastfoot
نمادین / مبتنی بر قاعده & شناخت به‌عنوان دست‌کاری نمادها و قواعد تولید
\lr{(production rules)} & سیستم‌های تولید \lr{ACT-R} \\
پیوندگرا \lr{(connectionist)} & بازنمایی‌های توزیع‌شده، یادگیری با تغییر
وزن‌ها & شبکه‌های قاعده‌ی دلتا، شبکه‌های عمیق \\
تصادفی / پویا \lr{(stochastic} / \lr{dynamical)} & فرایندهای پرنویز که
در زمان پیش می‌روند & \lr{random walk}، \lr{drift-diffusion}،
\lr{Rescorla\lr{--}Wagner} \\
بیزی (احتمالاتی) & شناخت به‌عنوان استنباط احتمالاتی & مدل‌های مشاهده‌گر
بیزی \\
\end{longtable}
}

\begin{quote}
\textbf{مراقب باشید:} «مدل‌سازی بیزی» در این درس دو معنای متفاوت دارد.
(الف) \textbf{مدل‌های بیزی شناخت} ادعا می‌کنند که \emph{ذهن} استنباط بیزی
انجام می‌دهد (هفته‌ی 15). (ب) \textbf{تحلیل داده‌ی بیزی} روشی است که
\emph{ما پژوهشگران} با آن پارامترهای هر مدلی را، بیزی یا غیربیزی، برآورد
می‌کنیم (هفته‌های 9\lr{--}13). هفته‌های 9\lr{--}13 بیشتر درباره‌ی (ب) هستند.
\end{quote}

\subsubsection{\texorpdfstring{4. سه سطح تحلیل
\lr{Marr}}{4. سه سطح تحلیل Marr}}

\lr{David Marr} پیشنهاد کرد که یک سیستم پردازش اطلاعات را در سه سطح مکمل
تحلیل کنیم:

{\def\LTcaptype{none} % do not increment counter
\begin{longtable}[]{@{}
  >{\raggedright\arraybackslash}p{(\linewidth - 4\tabcolsep) * \real{0.3333}}
  >{\raggedright\arraybackslash}p{(\linewidth - 4\tabcolsep) * \real{0.3333}}
  >{\raggedright\arraybackslash}p{(\linewidth - 4\tabcolsep) * \real{0.3333}}@{}}
\toprule\noalign{}
\begin{minipage}[b]{\linewidth}\raggedright
سطح
\end{minipage} & \begin{minipage}[b]{\linewidth}\raggedright
سؤال
\end{minipage} & \begin{minipage}[b]{\linewidth}\raggedright
مثال تصمیم‌گیری
\end{minipage} \\
\midrule\noalign{}
\endhead
\bottomrule\noalign{}
\endlastfoot
\textbf{محاسباتی \lr{(computational)}} & چه مسئله‌ای حل می‌شود و چرا؟ &
انتخاب گزینه‌ی درست با برقراری تعادل بین سرعت و دقت \\
\textbf{الگوریتمی / بازنمایی \lr{(algorithmic)}} & چه بازنمایی‌ها و
رویه‌هایی آن را حل می‌کنند؟ & انباشت شواهد پرنویز تا رسیدن به یک مرز \\
\textbf{پیاده‌سازی \lr{(implementational)}} & به‌طور فیزیکی چگونه محقق
می‌شود؟ & نرخ شلیک نورون‌ها که تا یک آستانه بالا می‌رود \\
\end{longtable}
}

مدل‌های هنجاری معمولاً در سطح محاسباتی‌اند، مدل‌های فرایندی در سطح
الگوریتمی، و مدل‌های علوم اعصاب محاسباتی در سطح پیاده‌سازی. بیشتر مدل‌های
این درس \textbf{در سطح الگوریتمی} هستند.

\subsubsection{5. مدل یک ماشین تولید داده است}

مهم‌ترین ایده برای ادامه‌ی درس: \textbf{مدل یک توزیع احتمال روی داده‌های
ممکن تعریف می‌کند.}

\[
\mathbf{y} \sim p(\mathbf{y} \mid \boldsymbol{\theta}, \mathcal{M})
\]

این را این‌طور بخوانید: «تحت مدل \(\mathcal{M}\) با پارامترهای
\(\boldsymbol{\theta}\)، داده‌ی \(\mathbf{y}\) از توزیع \(p\) تولید
می‌شود.» بسیاری از مدل‌ها این را به یک هسته‌ی قطعی به‌علاوه‌ی نویز تقسیم
می‌کنند، مثلاً \(y_i = f(x_i; \boldsymbol{\theta}) + \varepsilon_i\)
(هفته‌ی 2). بعضی دیگر، مثل \lr{random walk}، از اول تا آخر تصادفی‌اند.

این دو جهت استدلال به ما می‌دهد:

{\def\LTcaptype{none} % do not increment counter
\begin{longtable}[]{@{}
  >{\raggedright\arraybackslash}p{(\linewidth - 8\tabcolsep) * \real{0.2000}}
  >{\raggedright\arraybackslash}p{(\linewidth - 8\tabcolsep) * \real{0.2000}}
  >{\raggedright\arraybackslash}p{(\linewidth - 8\tabcolsep) * \real{0.2000}}
  >{\raggedright\arraybackslash}p{(\linewidth - 8\tabcolsep) * \real{0.2000}}
  >{\raggedright\arraybackslash}p{(\linewidth - 8\tabcolsep) * \real{0.2000}}@{}}
\toprule\noalign{}
\begin{minipage}[b]{\linewidth}\raggedright
جهت
\end{minipage} & \begin{minipage}[b]{\linewidth}\raggedright
داده‌شده
\end{minipage} & \begin{minipage}[b]{\linewidth}\raggedright
خواسته‌شده
\end{minipage} & \begin{minipage}[b]{\linewidth}\raggedright
نام
\end{minipage} & \begin{minipage}[b]{\linewidth}\raggedright
هفته‌ها
\end{minipage} \\
\midrule\noalign{}
\endhead
\bottomrule\noalign{}
\endlastfoot
\textbf{مستقیم \lr{(forward)}} & \(\boldsymbol\theta\) & داده‌های ممکن
\(\mathbf{y}\) & شبیه‌سازی / پیش‌بینی & 1\lr{--}2 \\
\textbf{معکوس \lr{(inverse)}} & \(\mathbf{y}\) مشاهده‌شده &
\(\boldsymbol\theta\) محتمل & برآورد پارامتر / برازش & 3\lr{--}5، 9\lr{--}13 \\
\end{longtable}
}

مسئله‌ی معکوس سخت‌تر است و همان چیزی است که منظورمان از «برازش مدل روی
داده» است. چون همیشه می‌توان مدل را در جهت مستقیم اجرا کرد، می‌توانیم
روش‌های برازشمان را هم آزمون کنیم: از پارامترهای \textbf{معلوم}
\(\boldsymbol\theta^\ast\) داده شبیه‌سازی کنیم، برازش کنیم، و ببینیم آیا
\(\boldsymbol\theta^\ast\) را پس می‌گیریم (\emph{بازیابی پارامتر} یا
\lr{parameter recovery}، هفته‌ی 7).

\paragraph{خواندن ریاضی}

\(\mathbf{y} \sim p(\mathbf{y} \mid \boldsymbol{\theta},\allowbreak \mathcal{M})\)
این‌طور خوانده می‌شود: «داده‌ی \(\mathbf{y}\) از توزیع احتمال \(p\) تولید
می‌شود، به شرط مقادیر پارامتر \(\boldsymbol\theta\) و مدل
\(\mathcal{M}\).»

\begin{itemize}
\tightlist
\item
  \(\mathbf{y}\) (پررنگ) یک \emph{فهرست} از مشاهده‌هاست؛ مثلاً 200 پاسخ یک
  شرکت‌کننده.
\item
  علامت \(\sim\) یعنی «به‌طور تصادفی تولید می‌شود از». یادآوری می‌کند که یک
  مدل با پارامترهای ثابت هر بار داده‌ی \emph{متفاوتی} می‌دهد، همان‌طور که
  یک آدم هم در هر \lr{trial} دقیقاً یکسان پاسخ نمی‌دهد.
\item
  خط عمودی \(\mid\) یعنی «به شرطِ» یا «با فرضِ». هر چیزی سمت راستش ثابت
  است: معادله‌های مدل (\(\mathcal{M}\)) و عددهایی که در آن‌ها می‌گذاریم
  (\(\boldsymbol\theta\)).
\item
  \(\boldsymbol\theta\) (تتا، پررنگ) همه‌ی پارامترها را کنار هم می‌گذارد؛
  مثلاً برای \lr{random walk} داریم \(\boldsymbol\theta = (v, a, z)\).
\end{itemize}

\textbf{مثال.} برای سکه‌ای با \(P(\text{\lr{heads}}) = 0.6\)، مدل این است: «هر
پرتاب با احتمال \(\theta\) شیر می‌آید»، \(\boldsymbol\theta = 0.6\)، و
\(\mathbf{y}\) می‌تواند (شیر، خط، شیر، شیر، خط) باشد. دوباره اجرا کنید،
دنباله‌ی دیگری می‌گیرید. جهت \emph{مستقیم}: \(\theta\) را انتخاب کن و
\(\mathbf{y}\) را شبیه‌سازی کن. جهت \emph{معکوس} (برازش، هفته‌های 3 تا 5):
\(\mathbf{y}\) را ببین و یک \(\theta\) خوب پیدا کن.

\subsubsection{6. مسیر مدل‌سازی (نقشه‌ی راه ما)}

\Needspace{13\baselineskip}
% Vector rendering of the supplied week-01 Mermaid pipeline.
\begin{figure}[htbp]
\centering
\begin{tikzpicture}[>=stealth, every node/.style={draw=CMForest,rounded corners=2pt,fill=CMPale,
    text=CMInk,align=center,minimum width=3.2cm,minimum height=1.05cm,font=\small},
    every path/.style={draw=CMEmerald,thick}]
\ifCMpersian
  \node (a) at (0,0) {نظریه‌ی کلامی};
  \node (b) at (4.2,0) {مدل فرمال\\(هفته‌ی ۱\lr{--}۲)};
  \node (c) at (8.4,0) {شبیه‌سازی\\(هفته‌ی ۱\lr{--}۲)};
  \node (d) at (12.6,0) {برازش روی داده\\(هفته‌ی ۳\lr{--}۵، ۹\lr{--}۱۳)};
  \node (e) at (12.6,-2.0) {ارزیابی و مقایسه\\(هفته‌ی ۶\lr{--}۷، ۱۰)};
  \node (f) at (8.4,-2.0) {بازیابی پارامتر و مدل\\(هفته‌ی ۷)};
  \node (g) at (4.2,-2.0) {تفسیر\\(هر هفته)};
\else
  \node (a) at (0,0) {Verbal theory};
  \node (b) at (4.2,0) {Formal model\\(Weeks 1–2)};
  \node (c) at (8.4,0) {Simulate\\(Weeks 1–2)};
  \node (d) at (12.6,0) {Fit to data\\(Weeks 3–5, 9–13)};
  \node (e) at (12.6,-2.0) {Evaluate and compare\\(Weeks 6–7, 10)};
  \node (f) at (8.4,-2.0) {Recover parameters and models\\(Week 7)};
  \node (g) at (4.2,-2.0) {Interpret\\(every week)};
\fi
\path[->] (a) -- (b);
\path[->] (b) -- (c);
\path[->] (c) -- (d);
\path[->] (d) -- (e);
\path[->] (e) -- (f);
\path[->] (f) -- (g);
\path[->,dashed] (g) -- (b);
\end{tikzpicture}
\ifCMpersian\caption{چرخه‌ی مدل‌سازی شناختی، بازطراحی شده از نمودار هفته‌ی اول}\else\caption{The cognitive modeling pipeline from Week 1}\fi
\label{fig:modeling-pipeline}
\end{figure}


هر \lr{lab} در این درس نسخه‌ای از این چرخه را اجرا می‌کند. امروز روی
\textbf{فرمال‌سازی \ensuremath{\leftarrow} شبیه‌سازی \ensuremath{\leftarrow} تفسیر} تمرکز می‌کنیم، به‌علاوه‌ی پیش‌نمایش
کوچکی از \textbf{برازش}.

\subsubsection{\texorpdfstring{7. اولین مدل ما: \lr{random walk} تا
رسیدن به مرز}{7. اولین مدل ما: random walk تا رسیدن به مرز}}

\textbf{از کلامی به فرمال.} شواهد یک عدد واحد \(X\) است که از \(z\) شروع
می‌شود. در هر گام زمانی \textbf{یک واحد بالا} می‌رود (شواهد به نفع پاسخ
درست) یا \textbf{یک واحد پایین} (شواهد به نفع پاسخ غلط). وقتی \(X\) به
مرز بالا \(a\) (پاسخ درست) یا مرز پایین \(0\) (خطا) برسد، حرکت متوقف
می‌شود.

\[
X_0 = z, \qquad X_{n} = X_{n-1} + \delta_n, \qquad
\delta_n =
\begin{cases}
+1 & \text{\lr{with probability }} \tfrac{1+v}{2}\\[2pt]
-1 & \text{\lr{with probability }} \tfrac{1-v}{2}
\end{cases}
\]

\begin{itemize}
\tightlist
\item
  \(v \in (-1, 1)\) همان \textbf{\lr{drift}} است، یعنی میانگین گام
  \(\mathbb{E}[\delta_n] = v\). نماینده‌ی \textbf{کیفیت شواهد} است:
  محرک‌های آسان \(v\) بزرگ می‌دهند.
\item
  \(a\) \textbf{فاصله‌ی مرزها} \lr{(boundary separation)} است و نماینده‌ی
  \textbf{احتیاط}: فاصله‌ی بیشتر یعنی شواهد بیشتری لازم است.
\item
  \(z\) \textbf{نقطه‌ی شروع} است (\(0 < z < a\)). با \(z = a/2\)
  تصمیم‌گیرنده بی‌سوگیری است.
\item
  تعداد گام‌ها تا رسیدن به مرز، \(T\)، از این راه به زمان واکنش نگاشت
  می‌شود: \[
  \mathrm{RT} = t_0 + T\,\Delta t,
  \] که در آن \(t_0\) \textbf{\lr{non-decision time}} \lr{(non-decision
  time:} رمزگذاری به‌علاوه‌ی اجرای حرکتی) و \(\Delta t\) مدت هر گام است.
\end{itemize}

\textbf{پیش‌بینی‌های تحلیلی.} فرض کنید \(r = \frac{1-v}{1+v}\). برای
\(v \neq 0\)، احتمال پاسخ درست و میانگین تعداد گام‌ها برابرند با:

\[
P_c = \frac{1 - r^{z}}{1 - r^{a}}, \qquad
\mathbb{E}[T] = \frac{a\,P_c - z}{v}.
\]

همین دو معادله سؤال‌های بخش انگیزه را حل می‌کنند:

\begin{itemize}
\tightlist
\item
  افزایش \(a\) هم \(P_c\) \textbf{و هم} \(\mathbb{E}[T]\) را بالا می‌برد:
  یک \textbf{\lr{trade-off} سرعت\lr{--}دقت} \lr{(speed\lr{--}accuracy trade-off).}
\item
  افزایش \(v\) احتمال \(P_c\) را بالا می‌برد ولی \(\mathbb{E}[T]\) را
  \textbf{پایین} می‌آورد: محرک‌های آسان‌تر هم سریع‌ترند و هم دقیق‌تر.
\end{itemize}

پس این دو دست‌کاری ردپای متفاوتی در داده‌ها می‌گذارند. همین است که به ما
اجازه می‌دهد احتیاط و کیفیت شواهد را جداگانه \emph{اندازه بگیریم}.

\begin{quote}
\textbf{پیشرفته: فرمول‌ها از کجا می‌آیند.} \(P_c\) همان نتیجه‌ی کلاسیک
\emph{ورشکستگی قمارباز} \lr{(gambler's ruin)} است. فرض کنید \(h(x)\)
احتمال رسیدن به \(a\) پیش از \(0\) با شروع از \(x\) باشد. شرطی کردن روی
گام اول می‌دهد \(h(x) = \tfrac{1+v}{2}h(x+1) + \tfrac{1-v}{2}h(x-1)\) با
\(h(0)=0\) و \(h(a)=1\). حل این معادله‌ی تفاضلی می‌دهد
\(h(x) = (1-r^x)/(1-r^a)\). برای \(\mathbb{E}[T]\) از \emph{اتحاد
\lr{Wald}} استفاده کنید: موقعیت نهایی در
\(\mathbb{E}[X_T] = z + v\,\mathbb{E}[T]\) صدق می‌کند. چون \(X_T\) با
احتمال \(P_c\) برابر \(a\) و در غیر این صورت \(0\) است،
\(\mathbb{E}[X_T] = aP_c\)، و مرتب کردن معادله می‌دهد
\(\mathbb{E}[T] = (aP_c - z)/v\). وقتی اندازه‌ی گام و \(\Delta t\) با هم
کوچک شوند، \lr{random walk} به \textbf{مدل \lr{drift-diffusion}} (هفته‌ی
14) تبدیل می‌شود.
\end{quote}

\begin{quote}
\textbf{پیشرفته: ویژگی‌ای که نقطه‌ضعف از آب درمی‌آید.} با شروع بی‌سوگیری
(\(z = a/2\))، \lr{random walk} ساده پیش‌بینی می‌کند که پاسخ‌های درست و غلط
\textbf{توزیع زمان واکنش یکسانی} دارند. در داده‌های واقعی، وقتی بر دقت
تأکید می‌شود خطاها اغلب \emph{کندتر} از پاسخ‌های درست‌اند و گاهی تحت تأکید
بر سرعت \emph{سریع‌تر}. مدل ساده هیچ‌کدام از این الگوها را نمی‌تواند تولید
کند. اضافه کردن تغییرپذیری \lr{drift} یا نقطه‌ی شروع از \lr{trial}ی به
\lr{trial} دیگر می‌تواند (تمرین 5 در \lr{lab).} پیدا کردن جایی که مدل
شکست می‌خورد دقیقاً همان راه بهتر شدن مدل‌هاست.
\end{quote}

\paragraph{خواندن ریاضی}

\textbf{قاعده‌ی به‌روزرسانی.} \(X_n = X_{n-1} + \delta_n\) یعنی:
\emph{شواهد الان} = \emph{شواهد یک قدم قبل} + \emph{یک قدم جدید}. قدم
\(\delta_n\) برابر \(+1\) یا \(-1\) است. احتمال \(+1\) برابر
\(\tfrac{1+v}{2}\) است. با \(v = 0\) این احتمال \(0.5\) است (سکه‌ی سالم:
هیچ اطلاعاتی نیست). با \(v = 0.2\) برابر \(0.6\) است، پس مسیر به‌طور
میانگین رو به بالا می‌رود. میانگین قدم برابر
\((+1)(0.6) + (-1)(0.4) = 0.2 = v\) است؛ به همین دلیل به \(v\) می‌گوییم
\lr{drift} (رانش).

\textbf{فرمول دقت.} \(P_c = \dfrac{1 - r^{z}}{1 - r^{a}}\) با
\(r = \dfrac{1-v}{1+v}\).

\begin{itemize}
\tightlist
\item
  \(r\) نسبت «احتمال قدم پایین / احتمال قدم بالا» است. با \(v = 0.2\)
  داریم \(r = 0.4/0.6 = 2/3\). هرچه \(r\) کوچک‌تر باشد، مسیر تقریباً همیشه
  بالا می‌رود.
\item
  وقتی \(z\) و \(a\) بزرگ باشند، \(r^{z}\) و \(r^{a}\) عددهای کوچکی
  می‌شوند و \(P_c\) به 1 نزدیک می‌شود. این همان ریاضیِ پشت جمله‌ی «احتیاط
  بیشتر (یعنی \(a\) بزرگ‌تر) دقت بیشتری می‌دهد» است.
\item
  لازم نیست فرمول را اثبات کنید (کادر پیشرفته را ببینید). کافی است عدد
  بگذارید و آن را با شبیه‌سازی مقایسه کنید، همان کاری که در \lr{lab}
  می‌کنیم.
\end{itemize}

\textbf{فرمول زمان.} \(\mathbb{E}[T] = \dfrac{a P_c - z}{v}\).
\(\mathbb{E}[\cdot]\) یعنی «میانگین روی تعداد زیادی \lr{trial».} صورت
کسر مسافتی است که مسیر به‌طور میانگین طی می‌کند (با احتمال \(P_c\) در
\(a\) و در غیر این صورت در \(0\) تمام می‌شود و از \(z\) شروع کرده بود).
تقسیم بر سرعت میانگین \(v\) تعداد میانگین قدم‌ها را می‌دهد؛ درست مثل
\emph{زمان = مسافت / سرعت}.

\textbf{زمان واکنش.} \(\mathrm{RT} = t_0 + T\,\Delta t\): یک بخش ثابت
\(t_0\) (دیدن محرک و فشار دادن کلید) به‌علاوه‌ی زمان تصمیم، یعنی تعداد
قدم‌ها ضرب در مدت هر قدم.

\paragraph{برای یک مهندس}

شما این مدل را با نام دیگری می‌شناسید.

\begin{itemize}
\tightlist
\item
  \textbf{آشکارسازی ترتیبی \lr{(sequential detection).}} \lr{random
  walk} نسخه‌ی گسسته‌ی \textbf{آزمون نسبت احتمال ترتیبی \lr{Wald (SPRT)}}
  است: نسبت لگاریتم درست‌نمایی مشاهده‌ها را جمع کن و به محض اینکه از یک
  آستانه‌ی بالا یا پایین گذشت، تصمیم بگیر. تحت بعضی فرض‌ها، \lr{SPRT} برای
  احتمال خطای مشخص، کمترین تعداد میانگین مشاهده را لازم دارد؛ یکی از
  دلایلی که ممکن است مغز از این راهبرد استفاده کند.
\item
  \textbf{\lr{drift} \ensuremath{\approx} نسبت سیگنال به نویز.} هر قدم «سیگنال + نویز» است.
  \lr{drift} \(v\) میانگین قدم است و تصادفی بودن قدم همان نویز است، پس
  \(v\) نقش \lr{SNR} هر نمونه را دارد. محرک واضح‌تر یعنی \lr{SNR} بالاتر.
\item
  \textbf{فاصله‌ی مرزها \ensuremath{\approx} آستانه‌ی تصمیم.} بالا بردن آستانه احتمال خطا را
  کم ولی تأخیر آشکارسازی را زیاد می‌کند: همان بده‌بستان کلاسیک بین نرخ خطا
  و تأخیر در آشکارسازی ترتیبی. در روان‌شناسی به این \textbf{بده‌بستان سرعت
  و دقت} می‌گویند.
\item
  \textbf{نقطه‌ی شروع \ensuremath{\approx} \lr{prior.}} شروع از نزدیکی یکی از مرزها مثل این
  است که نسبت لگاریتم درست‌نمایی اولیه صفر نباشد؛ یعنی یک \lr{prior} به
  نفع یکی از فرضیه‌ها.
\item
  \textbf{\lr{non-decision time} \ensuremath{\approx} تأخیر ثابت پردازش.} تأخیر
  حسگر/رمزگذاری به‌علاوه‌ی تأخیر عملگر (حرکت دست) که به هر پاسخ اضافه
  می‌شود.
\item
  \textbf{چیزی که برای یک مهندس تازه است:} در مهندسی ما آشکارساز را
  \emph{طراحی} می‌کنیم. اینجا یک آشکارساز انسانی را \emph{مشاهده} می‌کنیم
  و سعی می‌کنیم پارامترهایش \lr{(SNR}، آستانه) را از روی تصمیم‌ها و زمان
  تصمیم‌هایش \textbf{استنباط} کنیم. این مسئله‌ی معکوس موضوع کل این درس
  است.
\end{itemize}

\subsubsection{8. نمادگذاری کل درس}

این جدول در \textbf{همه‌ی} هفته‌ها استفاده می‌شود. وقتی هفته‌ای نمادی را
معرفی می‌کند، همان معنای این جدول را دارد.

\textbf{داده و طرح آزمایش}

{\def\LTcaptype{none} % do not increment counter
\begin{longtable}[]{@{}
  >{\raggedright\arraybackslash}p{(\linewidth - 4\tabcolsep) * \real{0.3333}}
  >{\raggedright\arraybackslash}p{(\linewidth - 4\tabcolsep) * \real{0.3333}}
  >{\raggedright\arraybackslash}p{(\linewidth - 4\tabcolsep) * \real{0.3333}}@{}}
\toprule\noalign{}
\begin{minipage}[b]{\linewidth}\raggedright
نماد
\end{minipage} & \begin{minipage}[b]{\linewidth}\raggedright
معنا
\end{minipage} & \begin{minipage}[b]{\linewidth}\raggedright
اولین استفاده
\end{minipage} \\
\midrule\noalign{}
\endhead
\bottomrule\noalign{}
\endlastfoot
\(y_i\)، \(\mathbf{y} = (y_1,\dots,y_N)\) & مشاهده‌ی \(i\)؛ بردار همه‌ی
مشاهدات & هفته‌ی 1 \\
\(N\) & تعداد مشاهدات \lr{(trial}ها) & هفته‌ی 1 \\
\(x_i\) & شرط آزمایشی / متغیر مستقل مشاهده‌ی \(i\) & هفته‌ی 2 \\
\(j = 1,\dots,J\) & اندیس شرکت‌کننده؛ تعداد شرکت‌کنندگان & هفته‌ی 5 \\
\end{longtable}
}

\textbf{مدل‌ها و پارامترها}

{\def\LTcaptype{none} % do not increment counter
\begin{longtable}[]{@{}
  >{\raggedright\arraybackslash}p{(\linewidth - 4\tabcolsep) * \real{0.3333}}
  >{\raggedright\arraybackslash}p{(\linewidth - 4\tabcolsep) * \real{0.3333}}
  >{\raggedright\arraybackslash}p{(\linewidth - 4\tabcolsep) * \real{0.3333}}@{}}
\toprule\noalign{}
\begin{minipage}[b]{\linewidth}\raggedright
نماد
\end{minipage} & \begin{minipage}[b]{\linewidth}\raggedright
معنا
\end{minipage} & \begin{minipage}[b]{\linewidth}\raggedright
اولین استفاده
\end{minipage} \\
\midrule\noalign{}
\endhead
\bottomrule\noalign{}
\endlastfoot
\(\mathcal{M}\)، \(\mathcal{M}_1, \mathcal{M}_2\) & یک مدل؛ مدل‌های رقیب
& هفته‌ی 1 \\
\(\boldsymbol\theta\) & بردار پارامترهای آزاد & هفته‌ی 1 \\
\(k\) & تعداد پارامترهای آزاد & هفته‌ی 1 \\
\(\boldsymbol\theta^\ast\) & مقادیر واقعی (تولیدکننده‌ی داده) پارامترها
در شبیه‌سازی & هفته‌ی 1 \\
\(\hat{\boldsymbol\theta}\) & برآورد \(\boldsymbol\theta\) & هفته‌ی 3 \\
\(f(x;\boldsymbol\theta)\) & بخش قطعی پیش‌بینی مدل & هفته‌ی 2 \\
\(\hat y_i\) & پیش‌بینی مدل برای مشاهده‌ی \(i\) & هفته‌ی 3 \\
\(\varepsilon_i\) & جمله‌ی نویز (خطا) & هفته‌ی 2 \\
\end{longtable}
}

\textbf{احتمال}

{\def\LTcaptype{none} % do not increment counter
\begin{longtable}[]{@{}
  >{\raggedright\arraybackslash}p{(\linewidth - 2\tabcolsep) * \real{0.5000}}
  >{\raggedright\arraybackslash}p{(\linewidth - 2\tabcolsep) * \real{0.5000}}@{}}
\toprule\noalign{}
\begin{minipage}[b]{\linewidth}\raggedright
نماد
\end{minipage} & \begin{minipage}[b]{\linewidth}\raggedright
معنا
\end{minipage} \\
\midrule\noalign{}
\endhead
\bottomrule\noalign{}
\endlastfoot
\(P(\cdot)\) & احتمال یک رویداد \\
\(p(\cdot)\) & تابع جرم یا چگالی احتمال \\
\(y \sim \mathcal{D}\) & «\(y\) توزیع \(\mathcal{D}\) دارد» \\
\(\mathcal{N}(\mu, \sigma)\) & توزیع نرمال با میانگین \(\mu\) و
\textbf{انحراف معیار} \(\sigma\) (مثل \lr{scipy} و \lr{PyMC)} \\
\(\mathrm{Bernoulli}(\theta)\)، \(\mathrm{Binomial}(n, \theta)\) &
توزیع‌های برنولی و دوجمله‌ای با احتمال موفقیت \(\theta\) \\
\(\mathrm{Beta}(a_0, b_0)\)، \(\mathrm{Uniform}(l, u)\) & توزیع‌های
\lr{Beta} و یکنواخت \\
\(\Phi\)، \(\Phi^{-1}\) & تابع توزیع تجمعی \lr{(CDF)} نرمال استاندارد و
معکوس آن \\
\(\mathbb{E}[\cdot]\)، \(\mathrm{Var}[\cdot]\) & امید ریاضی و واریانس \\
\end{longtable}
}

\textbf{برازش، ارزیابی و استنباط بیزی}

{\def\LTcaptype{none} % do not increment counter
\begin{longtable}[]{@{}
  >{\raggedright\arraybackslash}p{(\linewidth - 4\tabcolsep) * \real{0.3333}}
  >{\raggedright\arraybackslash}p{(\linewidth - 4\tabcolsep) * \real{0.3333}}
  >{\raggedright\arraybackslash}p{(\linewidth - 4\tabcolsep) * \real{0.3333}}@{}}
\toprule\noalign{}
\begin{minipage}[b]{\linewidth}\raggedright
نماد
\end{minipage} & \begin{minipage}[b]{\linewidth}\raggedright
معنا
\end{minipage} & \begin{minipage}[b]{\linewidth}\raggedright
اولین استفاده
\end{minipage} \\
\midrule\noalign{}
\endhead
\bottomrule\noalign{}
\endlastfoot
\lr{SSE}، \lr{MSE}، \lr{RMSD} & مجموع مربعات خطا، میانگین مربعات خطا،
ریشه‌ی میانگین مربع انحراف & هفته‌ی 3 \\
\(L(\boldsymbol\theta) = p(\mathbf{y}\mid\boldsymbol\theta,\mathcal{M})\)
& درست‌نمایی \lr{(likelihood)} & هفته‌ی 4 \\
\(\ell(\boldsymbol\theta) = \ln L(\boldsymbol\theta)\)؛ \lr{NLL}
\(= -\ell\) & لگاریتم درست‌نمایی \lr{(log-likelihood)}؛ منفی لگاریتم
درست‌نمایی & هفته‌ی 4 \\
\(K\) & تعداد بخش‌ها \lr{(folds)} در \lr{cross-validation} (با \(k\) یکی
نیست!) & هفته‌ی 6 \\
\lr{AIC}، \lr{BIC} & معیارهای اطلاعاتی آکائیکه و بیزی & هفته‌ی 7 \\
\(p(\boldsymbol\theta)\)، \(p(\boldsymbol\theta\mid\mathbf{y})\) & پیشین
\lr{(prior)}؛ پسین \lr{(posterior)} & هفته‌ی 9 \\
\(p(\mathbf{y}\mid\mathcal{M})\) & \lr{marginal likelihood} (شواهد یا
\lr{evidence)} & هفته‌ی 10 \\
\(\mathrm{BF}_{12}\) & \lr{Bayes factor} برای \(\mathcal{M}_1\) در برابر
\(\mathcal{M}_2\) & هفته‌ی 10 \\
\(\boldsymbol\theta^{(s)}\)، \(s = 1,\dots,S\) & نمونه‌ی \(s\)-ام
\lr{MCMC}؛ تعداد نمونه‌ها & هفته‌ی 11 \\
\(C\) & تعداد زنجیره‌های \lr{MCMC} & هفته‌ی 12 \\
\(\hat R\)، \lr{ESS} & ضریب کاهش مقیاس بالقوه \lr{(R-hat)}؛ اندازه‌ی مؤثر
نمونه & هفته‌ی 12 \\
\end{longtable}
}

\textbf{نمادهای مخصوص هر مدل}

{\def\LTcaptype{none} % do not increment counter
\begin{longtable}[]{@{}
  >{\raggedright\arraybackslash}p{(\linewidth - 4\tabcolsep) * \real{0.3333}}
  >{\raggedright\arraybackslash}p{(\linewidth - 4\tabcolsep) * \real{0.3333}}
  >{\raggedright\arraybackslash}p{(\linewidth - 4\tabcolsep) * \real{0.3333}}@{}}
\toprule\noalign{}
\begin{minipage}[b]{\linewidth}\raggedright
مدل
\end{minipage} & \begin{minipage}[b]{\linewidth}\raggedright
نمادها
\end{minipage} & \begin{minipage}[b]{\linewidth}\raggedright
اولین استفاده
\end{minipage} \\
\midrule\noalign{}
\endhead
\bottomrule\noalign{}
\endlastfoot
\lr{random walk} / \lr{drift-diffusion} & \(v\) \lr{drift rate}؛ \(a\)
فاصله‌ی مرزها؛ \(z\) نقطه‌ی شروع؛ \(t_0\) \lr{non-decision time}؛
\(\Delta t\) مدت هر گام؛ \(T\) تعداد گام‌ها & هفته‌ی 1 و 14 \\
نظریه‌ی تشخیص علامت & \(H\) نرخ \lr{hit}؛ \(F\) نرخ \lr{false alarm}؛
\(d'\) حساسیت؛ \(c\) ملاک \lr{(criterion)} & هفته‌ی 2 \\
منحنی‌های فراموشی & \(t\) فاصله‌ی نگهداری \lr{(retention interval)}؛ \(A\)
سطح اولیه؛ \(\lambda\) نرخ زوال نمایی؛ \(d\) نمای زوال توانی & هفته‌ی
2 \\
\lr{GCM} (دسته‌بندی) & \(w_m\) وزن‌های توجه؛ \(c\) حساسیت & هفته‌ی 14 \\
یادگیری تقویتی & \(\alpha\) نرخ یادگیری؛ \(\beta\) \lr{inverse
temperature}؛ \(Q_t(o)\) ارزش گزینه‌ی \(o\) در \lr{trial} \(t\) & هفته‌ی
15 \\
\end{longtable}
}

\textbf{نمادهای دومعنایی (چون در منابع رایج‌اند نگه داشته شده‌اند):} \(c\)
ملاک \lr{SDT} است، جز در هفته‌ی 14 که حساسیت \lr{GCM} است. \(t\) در
مدل‌های فراموشی فاصله‌ی نگهداری و در یادگیری تقویتی اندیس \lr{trial} است.
\(\alpha\) فقط به معنای نرخ یادگیری در یادگیری تقویتی است. هیچ‌وقت
\(\alpha\) را برای سطح معناداری به کار نمی‌بریم.

\paragraph{برای مدرس}

\begin{itemize}
\tightlist
\item
  \textbf{با سؤال‌ها شروع کنید، نه با مدل.} نظریه‌ی کلامی را بلند بخوانید
  و از دانشجوها بخواهید به چهار سؤال بخش مقدمه جواب دهند. با هم اختلاف
  نظر پیدا می‌کنند. همین اختلاف، دلیل وجود مدل‌های فرمال است.
\item
  \textbf{\lr{random walk} را اجرا کنید.} روی تخته یک خط عمودی بکشید، با
  یک خط بالا («درست») و یک خط پایین («خطا»). یک دانشجو تاس بیندازد: 1 تا
  4 یعنی بالا، 5 و 6 یعنی پایین (\(v \approx 0.33\)). قدم‌ها را تا رسیدن
  به یکی از خط‌ها بشمارید. سه بار تکرار کنید و سه \lr{trial} با زمان‌های
  واکنش متفاوت دارید. بعد خط‌ها را از هم دورتر کنید و دوباره امتحان کنید.
\item
  \textbf{در این هفته ریاضی را سبک نگه دارید.} دانشجوها فقط باید
  فرمول‌های \(P_c\) و \(\mathbb{E}[T]\) را \emph{استفاده} کنند، نه اثبات.
  در \lr{lab} نشان دهید که عددها با شبیه‌سازی جور درمی‌آیند؛ پیام اصلی
  همین است (شبیه‌سازی و تحلیل با هم توافق دارند).
\item
  \textbf{سؤال‌های رایج دانشجوها.}

  \begin{itemize}
  \tightlist
  \item
    \emph{«آیا شواهد واقعاً قدم‌های گسسته است؟»} نه، این یک ساده‌سازی است.
    هفته‌ی 14 نسخه‌ی پیوسته \lr{(drift-diffusion)} را نشان می‌دهد که
    پیش‌بینی‌های مشابهی دارد.
  \item
    \emph{«چرا میانگین زمان واکنش را با \lr{t-test} مقایسه نکنیم؟»}
    \lr{t-test} می‌گوید گروه‌ها \emph{فرق دارند}؛ مدل می‌گوید \emph{چطور}
    فرق دارند (احتیاط در برابر کیفیت شواهد).
  \item
    \emph{«کدام مدل درست است؟»} هیچ‌کدام. همه‌ی مدل‌ها ساده‌سازی‌اند؛ می‌پرسیم
    کدام مفید است (جمله‌ی \lr{Box).}
  \end{itemize}
\item
  \textbf{مدیریت زمان.} بخش‌های 3 و 4 (انواع مدل، \lr{Marr)} را در صورت
  نیاز می‌شود به 10 دقیقه کوتاه کرد. بخش 7 را حذف نکنید: کل \lr{lab} به
  آن وابسته است.
\end{itemize}

\begin{center}\rule{0.5\linewidth}{0.5pt}\end{center}

\subsection{مثال حل‌شده}

\textbf{شرایط.} یک تصمیم‌گیرنده‌ی بی‌سوگیری با \lr{drift} \(v = 0.2\)،
فاصله‌ی مرزها \(a = 6\) و نقطه‌ی شروع \(z = 3\). \lr{non-decision time}
\(t_0 = 300\) \lr{ms}؛ هر گام \(\Delta t = 20\) \lr{ms} طول می‌کشد.

\textbf{قدم 1: احتمال گام‌ها.} \(P(\delta = +1) = (1 + 0.2)/2 = 0.6\) و
\(P(\delta = -1) = 0.4\).

\textbf{قدم 2: شبیه‌سازی دستی یک \lr{trial.}} اعداد تصادفی یکنواخت
\(u_n\) بکشید. اگر \(u_n < 0.6\) بالا بروید، وگرنه پایین.

{\def\LTcaptype{none} % do not increment counter
\begin{longtable}[]{@{}llll@{}}
\toprule\noalign{}
گام \(n\) & \(u_n\) & \(\delta_n\) & \(X_n\) \\
\midrule\noalign{}
\endhead
\bottomrule\noalign{}
\endlastfoot
0 & \lr{---} & \lr{---} & 3 \\
1 & 0.31 & +1 & 4 \\
2 & 0.72 & \ensuremath{-}1 & 3 \\
3 & 0.55 & +1 & 4 \\
4 & 0.12 & +1 & 5 \\
5 & 0.48 & +1 & \textbf{6 = \lr{a}} \ensuremath{\leftarrow} توقف \\
\end{longtable}
}

مسیر به مرز بالا می‌رسد: یک پاسخ \textbf{درست} بعد از \(T = 5\) گام، پس
\(\mathrm{RT} = 300 + 5 \times 20 = 400\) \lr{ms.}

\textbf{قدم 3: دقت تحلیلی.} \(r = (1-0.2)/(1+0.2) = 2/3\)، پس
\(r^3 = 8/27 \approx 0.296\) و \(r^6 = 64/729 \approx 0.088\).

\[
P_c = \frac{1 - 8/27}{1 - 64/729} = \frac{1}{1 + 8/27} = \frac{27}{35} \approx 0.771.
\]

(وقتی \(z = a/2\)، فرمول به \(P_c = 1/(1 + r^{a/2})\) ساده می‌شود.)

\textbf{قدم 4: میانگین تحلیلی زمان واکنش.}

\[
\mathbb{E}[T] = \frac{6 \times 0.771 - 3}{0.2} = \frac{57/35}{0.2} \approx 8.14 \text{\lr{ steps}},
\qquad
\mathbb{E}[\mathrm{RT}] = 300 + 8.14 \times 20 \approx 462.9 \text{\lr{ ms}}.
\]

\textbf{قدم 5: دو دست‌کاری، دو ردپا.}

{\def\LTcaptype{none} % do not increment counter
\begin{longtable}[]{@{}
  >{\raggedright\arraybackslash}p{(\linewidth - 10\tabcolsep) * \real{0.1667}}
  >{\raggedright\arraybackslash}p{(\linewidth - 10\tabcolsep) * \real{0.1667}}
  >{\raggedright\arraybackslash}p{(\linewidth - 10\tabcolsep) * \real{0.1667}}
  >{\raggedright\arraybackslash}p{(\linewidth - 10\tabcolsep) * \real{0.1667}}
  >{\raggedright\arraybackslash}p{(\linewidth - 10\tabcolsep) * \real{0.1667}}
  >{\raggedright\arraybackslash}p{(\linewidth - 10\tabcolsep) * \real{0.1667}}@{}}
\toprule\noalign{}
\begin{minipage}[b]{\linewidth}\raggedright
شرط
\end{minipage} & \begin{minipage}[b]{\linewidth}\raggedright
\(v\)
\end{minipage} & \begin{minipage}[b]{\linewidth}\raggedright
\(a\) (\(z = a/2\))
\end{minipage} & \begin{minipage}[b]{\linewidth}\raggedright
\(P_c\)
\end{minipage} & \begin{minipage}[b]{\linewidth}\raggedright
\(\mathbb{E}[T]\)
\end{minipage} & \begin{minipage}[b]{\linewidth}\raggedright
\(\mathbb{E}[\mathrm{RT}]\) \lr{(ms)}
\end{minipage} \\
\midrule\noalign{}
\endhead
\bottomrule\noalign{}
\endlastfoot
پایه & 0.2 & 6 & 0.771 & 8.14 & 462.9 \\
«دقیق باشید!» (احتیاط بیشتر) & 0.2 & 10 & 0.884 & 19.18 & 683.6 \\
محرک آسان‌تر (شواهد بهتر) & 0.3 & 6 & 0.865 & 7.30 & 446.0 \\
\end{longtable}
}

\textbf{تفسیر.} دستور «دقیق باشید» حدود 11 درصد دقت بیشتر می‌خرد، به قیمت
حدود 220 \lr{ms.} محرک آسان‌تر افزایش دقت \emph{مشابهی} می‌دهد، در حالی که
پاسخ‌ها را کمی \emph{سریع‌تر} هم می‌کند. پژوهشگری که فقط به دقت نگاه کند
نمی‌تواند این دو وضعیت را از هم تشخیص دهد. مدل، با استفاده‌ی هم‌زمان از دقت
\emph{و} زمان واکنش، آن‌ها را به دو علت روان‌شناختی متفاوت نسبت می‌دهد
(احتیاط در برابر کیفیت شواهد).

\begin{center}\rule{0.5\linewidth}{0.5pt}\end{center}

\subsection{\texorpdfstring{\lr{lab} (2 ساعت)}{lab (2 ساعت)}}

\subsubsection{اهداف}

\begin{enumerate}
\def\labelenumi{\arabic{enumi}.}
\tightlist
\item
  نصب یک محیط نرم‌افزاری سالم روی لپ‌تاپ خودتان، \emph{از جمله} ابزارهای
  بیزی که در هفته‌ی 13 لازم‌اند.
\item
  تمرین کار پایه با \lr{numpy/matplotlib:} اعداد تصادفی، برداری‌سازی
  \lr{(vectorization)}، رسم نمودار.
\item
  پیاده‌سازی مدل \lr{random walk}، شبیه‌سازی یک آزمایش، و \textbf{وارسی}
  شبیه‌سازی با فرمول‌های تحلیلی.
\item
  استفاده از مدل برای \textbf{تفسیر} \lr{trade-off} سرعت\lr{--}دقت.
\item
  اولین تجربه‌ی \textbf{برازش}: بازیابی پارامترهای پنهان یک «شرکت‌کننده‌ی
  مرموز»، \textbf{ارزیابی} برازش و \textbf{تفسیر} نتیجه.
\end{enumerate}

\textbf{زمان‌بندی:} بخش \lr{A} 35 دقیقه · بخش \lr{B} 15 دقیقه · بخش
\lr{C} 30 دقیقه · بخش \lr{D} 15 دقیقه · بخش \lr{E} 25 دقیقه. تمرین‌هایی
که در \lr{lab} تمام نشوند در خانه تمام می‌شوند.

\subsubsection{\texorpdfstring{بخش \lr{A:} نصب محیط نرم‌افزاری (35
دقیقه)}{بخش A: نصب محیط نرم‌افزاری (35 دقیقه)}}

از \textbf{\lr{Miniforge}} استفاده می‌کنیم، یک نصب‌کننده‌ی رایگان برای
\lr{\nolinkurl{conda}} که از کانال \lr{\nolinkurl{conda-forge}} استفاده
می‌کند. این مطمئن‌ترین راه نصب \lr{PyMC} روی \lr{Windows}، \lr{macOS} و
\lr{Linux} است.

\begin{enumerate}
\def\labelenumi{\arabic{enumi}.}
\tightlist
\item
  \lr{Miniforge} را از سایت \lr{conda-forge} دانلود و نصب کنید
  (سیستم‌عامل خود را انتخاب کنید؛ گزینه‌های پیش‌فرض را نگه دارید).
\item
  فایل زیر را با نام \lr{\nolinkurl{environment.yml}} در پوشه‌ی درس ذخیره
  کنید:
\end{enumerate}

\Needspace{18\baselineskip}
\begin{latin}
\begin{Verbatim}
# Cognitive Modeling course environment (same as environment.yml in the course repository)
name: cogmod
channels:
  - conda-forge
dependencies:
  - python=3.12
  - numpy>=1.26
  - scipy>=1.11
  - pandas>=2.1
  - matplotlib>=3.8
  - jupyterlab
  - pymc>=6.3    # Bayesian modeling (needed from Week 13)
  - arviz>=1.3   # MCMC diagnostics
\end{Verbatim}
\end{latin}

\begin{enumerate}
\def\labelenumi{\arabic{enumi}.}
\setcounter{enumi}{2}
\tightlist
\item
  \textbf{\lr{Miniforge Prompt}} (در \lr{Windows)} یا یک \lr{terminal}
  (در \lr{macOS/Linux)} را باز کنید، به پوشه‌ی درس بروید و محیط را
  بسازید. بار اول 5 تا 15 دقیقه طول می‌کشد.
\end{enumerate}

\Needspace{6\baselineskip}
\begin{latin}
\begin{Verbatim}
conda env create -f environment.yml
\end{Verbatim}
\end{latin}

\Needspace{6\baselineskip}
\begin{latin}
\begin{Verbatim}
conda activate cogmod
\end{Verbatim}
\end{latin}

\begin{enumerate}
\def\labelenumi{\arabic{enumi}.}
\setcounter{enumi}{3}
\tightlist
\item
  این اسکریپت را با نام \lr{\nolinkurl{check_install.py}} ذخیره کنید و
  با \lr{\nolinkurl{python check_install.py}} اجرا کنید. با افزودن
  \lr{\nolinkurl{--sample}} یک مدل کوچک \lr{PyMC} هم اجرا می‌شود؛ اجرای
  اول حدود یک دقیقه طول می‌کشد چون \lr{PyMC} کامپایل می‌کند.
\end{enumerate}

\begin{quote}
\textbf{\lr{Code} · کد:} \lr{\nolinkurl{code/check_install.py}}
\lr{(save and run it as a separate script} · به‌صورت یک اسکریپت جدا اجرا
شود)
\end{quote}

\begin{enumerate}
\def\labelenumi{\arabic{enumi}.}
\setcounter{enumi}{4}
\tightlist
\item
  \lr{JupyterLab} را با \lr{\nolinkurl{jupyter lab}} اجرا کنید و یک
  \lr{notebook} با نام \lr{\nolinkurl{week01_lab.ipynb}} بسازید. بقیه‌ی
  کدها را در همین \lr{notebook}، هر بلوک در یک \lr{cell}، اجرا کنید.
\end{enumerate}

\begin{quote}
\textbf{اگر اسکریپت هشداری مثل \lr{«g++ not available»} نشان داد:}
\lr{PyMC} باز هم کار می‌کند ولی خیلی کندتر. معمولاً یعنی \lr{PyMC} با
\lr{\nolinkurl{pip}} نصب شده، نه از محیط \lr{conda-forge} بالا. نصب
\lr{conda-forge} معمولاً \lr{compiler} مورد نیاز \lr{PyMC} را هم می‌آورد
{[}پیش از هفته‌ی 1 روی یک کامپیوتر \lr{Windows} در \lr{lab} بررسی شود{]}.
محیط را دوباره با \lr{\nolinkurl{conda env create -f environment.yml}}
بسازید.
\end{quote}

\begin{quote}
\textbf{اگر نصب شکست خورد:} \lr{lab} امروز را در \textbf{\lr{Google
Colab}} انجام دهید \lr{(numpy}، \lr{scipy}، \lr{pandas} و
\lr{matplotlib} در آن از قبل نصب‌اند) و نصب محلی را پیش از هفته‌ی 3 با کمک
مدرس درست کنید. نصب \lr{PyMC} باید تا هفته‌ی 12 کار کند.
\end{quote}

\subsubsection{\texorpdfstring{بخش \lr{B:} گرم کردن با \lr{Python} (15
دقیقه)}{بخش B: گرم کردن با Python (15 دقیقه)}}

\begin{quote}
\textbf{\lr{Code} · کد:} \lr{\nolinkurl{code/week01.py}}, \lr{cell} 1 ·
سلول 1
\end{quote}

\textbf{خروجی مورد انتظار:} نسبتی نزدیک به 0.6 (با این \lr{seed} برابر
0.574) و منحنی نسبت تجمعی که اول نوسان دارد و بعد نزدیک خط‌چین آرام
می‌گیرد.

\subsubsection{\texorpdfstring{بخش \lr{C:} پیاده‌سازی، شبیه‌سازی و وارسی
\lr{random walk} (30
دقیقه)}{بخش C: پیاده‌سازی، شبیه‌سازی و وارسی random walk (30 دقیقه)}}

\begin{quote}
\textbf{\lr{Code} · کد:} \lr{\nolinkurl{code/week01.py}}, \lr{cell} 2 ·
سلول 2
\end{quote}

اول ببینید مدل چه می‌کند: چند مسیر شواهد را رسم کنید.

\begin{quote}
\textbf{\lr{Code} · کد:} \lr{\nolinkurl{code/week01.py}}, \lr{cell} 3 ·
سلول 3
\end{quote}

حالا یک آزمایش کامل را شبیه‌سازی کنید و نتیجه را با پیش‌بینی‌های تحلیلی
مقایسه کنید. \textbf{همیشه شبیه‌سازی را با چیزی که جوابش را می‌دانید
مقایسه کنید}؛ این عادت در این ترم بارها نجاتتان می‌دهد.

\begin{quote}
\textbf{\lr{Code} · کد:} \lr{\nolinkurl{code/week01.py}}, \lr{cell} 4 ·
سلول 4
\end{quote}

\begin{quote}
\textbf{\lr{Code} · کد:} \lr{\nolinkurl{code/week01.py}}, \lr{cell} 5 ·
سلول 5
\end{quote}

\textbf{خروجی مورد انتظار \lr{(seed} 42):}

\Needspace{14\baselineskip}
\begin{latin}
\begin{Verbatim}
Accuracy   simulated 0.764 | analytic 0.771
Mean steps simulated 8.27  | analytic 8.14
Mean RT    simulated 465.3 ms | analytic 462.9 ms

Mean RT by response (1 = correct, 0 = error):
        count   mean
choice
0        1180  466.1
1        3820  465.1
\end{Verbatim}
\end{latin}

\textbf{به این‌ها دقت کنید:} - شبیه‌سازی در حد نویز نمونه‌گیری با مقادیر
تحلیلی جور است. - توزیع‌های زمان واکنش \textbf{چوله به راست} هستند (یک
دنباله‌ی بلند از پاسخ‌های کند)؛ این شکل معمول توزیع‌های واقعی زمان واکنش
است. - زمان واکنش پاسخ‌های درست و غلط تقریباً یکسان است (کادر «پیشرفته» در
بخش 7 جزوه را ببینید).

\subsubsection{\texorpdfstring{بخش \lr{D:} تفسیر \lr{trade-off}
سرعت\lr{--}دقت (15 دقیقه)}{بخش D: تفسیر trade-off سرعت\lr{--}دقت (15 دقیقه)}}

\begin{quote}
\textbf{\lr{Code} · کد:} \lr{\nolinkurl{code/week01.py}}, \lr{cell} 6 ·
سلول 6
\end{quote}

\textbf{خروجی مورد انتظار:} سه منحنی که از پایین سمت چپ به سمت 1.0 بالا
می‌روند. حرکت روی یک منحنی (افزایش \(a\)) سرعت را با دقت معاوضه می‌کند.
حرکت \emph{بین} منحنی‌ها (افزایش \(v\)) دقت را با زمان واکنش برابر یا
کمتر بالا می‌برد. جدول مثال حل‌شده را بازتولید می‌کند: مثلاً
\(v = 0.2, a = 10\) دقت 0.884 و میانگین زمان واکنش 683.6 \lr{ms} می‌دهد.

\textbf{با بغل‌دستی‌تان بحث کنید:} پژوهشگری می‌بیند که یک دارو شرکت‌کننده‌ها
را کندتر \emph{و} دقیق‌تر می‌کند. به کدام پارامتر مشکوک می‌شوید؟ اگر دارو
آن‌ها را کندتر و \emph{کم‌دقت‌تر} می‌کرد چه؟

\subsubsection{\texorpdfstring{بخش \lr{E:} اولین برازش، ارزیابی و تفسیر
(25 دقیقه)}{بخش E: اولین برازش، ارزیابی و تفسیر (25 دقیقه)}}

مدرس فایل \lr{\nolinkurl{week01_mystery.csv}} را به شما می‌دهد که داده‌های
یک «شرکت‌کننده‌ی مرموز» است (1000 \lr{trial}، ستون‌های
\lr{\nolinkurl{choice}} و \lr{\nolinkurl{rt}}). اگر فایل را ندارید، کد
زیر یک مجموعه‌داده‌ی معادل با پارامترهایی که پشت یک \lr{seed} پنهان شده‌اند
می‌سازد. \textbf{دزدکی نگاه نکنید}: سعی کنید آن‌ها را بازیابی کنید.

فرض می‌کنیم \(t_0 = 300\) \lr{ms} و \(\Delta t = 20\) \lr{ms} معلوم‌اند و
\(z = a/2\). فقط \(v\) و \(a\) آزادند (\(k = 2\)).

\begin{quote}
\textbf{\lr{Code} · کد:} \lr{\nolinkurl{code/week01.py}}, \lr{cell} 7 ·
سلول 7
\end{quote}

\textbf{برازش با جستجوی شبکه‌ای \lr{(grid search).}} برای هر ترکیب \(v\)
و \(a\) روی یک شبکه، دقت و میانگین زمان واکنش پیش‌بینی‌شده را حساب کنید و
فاصله‌شان با مقادیر مشاهده‌شده را بسنجید. هر اختلاف بر حسب واحد خطای
استاندارد بیان می‌شود تا دقت و زمان واکنش مقیاس قابل مقایسه‌ای داشته
باشند. هفته‌ی 3 این ایده را به‌طور کامل توضیح می‌دهد.

\begin{quote}
\textbf{\lr{Code} · کد:} \lr{\nolinkurl{code/week01.py}}, \lr{cell} 8 ·
سلول 8
\end{quote}

\textbf{ارزیابی برازش.} فقط دو آماره‌ی خلاصه (دقت و میانگین زمان واکنش)
را برازش کردیم. آزمون قوی‌تر این است که ببینیم مدل چیزی را که برازش
\emph{نکرده‌ایم} هم بازتولید می‌کند یا نه: \textbf{شکل} توزیع زمان واکنش.
\lr{quantile}های زمان واکنش پاسخ‌های درست در داده را با \lr{quantile}های
یک شبیه‌سازی بزرگ با بهترین پارامترها مقایسه می‌کنیم.

\begin{quote}
\textbf{\lr{Code} · کد:} \lr{\nolinkurl{code/week01.py}}, \lr{cell} 9 ·
سلول 9
\end{quote}

\textbf{خروجی مورد انتظار (داده‌ی جایگزین). فقط بعد از اجرای برازش باز
کنید.}

\Needspace{14\baselineskip}
\begin{latin}
\begin{Verbatim}
Observed: accuracy = 0.836, mean RT = 746.4 ms (N = 1000)
Best fit: v = 0.15, a = 10 (pred. accuracy 0.819, pred. RT 725.7 ms, discrepancy 5.78)

   quantile  observed   model
0       0.1     440.0   440.0
1       0.3     520.0   520.0
2       0.5     640.0   640.0
3       0.7     800.0   800.0
4       0.9    1200.0  1160.0
\end{Verbatim}
\end{latin}

مقادیر پنهان \(v^\ast = 0.16\) و \(a^\ast = 10\) بودند. جستجوی شبکه‌ای
\(a\) را دقیقاً و \(v\) را با فاصله‌ی یک گام شبکه (0.15) بازیابی کرد. این
اختلاف کوچک باگ نیست: با 1000 \lr{trial}، دقت و میانگین زمان واکنش
مشاهده‌شده به خاطر نویز نمونه‌گیری با مقادیر واقعی فرق دارند، پس بهترین
پارامترها هم کمی فرق می‌کنند. (هفته‌ی 5 نشان می‌دهد چطور برای چنین
برآوردهایی بازه‌ی عدم قطعیت بسازیم.) زمان‌های واکنش مضرب 20 \lr{ms} هستند
و به همین دلیل بسیاری از \lr{quantile}ها دقیقاً برابرند. مدل شکل توزیع را
خوب بازتولید می‌کند و فقط دنباله‌ی کند را کمی کمتر برآورد می‌کند
\lr{(quantile} 0.9).

\textbf{سؤال‌های تفسیری (برای هرکدام 2 تا 3 جمله در \lr{notebook}
بنویسید):} 1. به زبان روان‌شناختی، \(v\) و \(a\) بازیابی‌شده درباره‌ی این
شرکت‌کننده چه می‌گویند؟ 2. برازش فقط از دو عدد داده استفاده کرد. چرا
مقایسه‌ی \lr{quantile}ها بررسی قانع‌کننده‌تری از خود مقدار ناهمخوانی
\lr{(discrepancy)} است؟ 3. اگر این شرکت‌کننده دوباره با دستور «تا جای
ممکن سریع باشید» آزمون شود، طبق مدل کدام پارامتر باید تغییر کند و کدام
ثابت بماند؟

\subsubsection{تمرین‌ها (به ترتیب سختی)}

\textbf{تمرین 1 \ensuremath{\star} (نویز نمونه‌گیری).} مدل پایه (\(v = 0.2\)، \(a = 6\))
را با \(N = 20\)، \(100\) و \(1000\) \lr{trial}، هرکدام با پنج \lr{seed}
مختلف، شبیه‌سازی کنید. جدولی از دقت‌های مشاهده‌شده بسازید. پراکندگی‌شان را
با خطای استاندارد \(\sqrt{P_c(1-P_c)/N}\) مقایسه کنید. \emph{انتظار:} با
\(N = 20\) دقت به‌راحتی بین حدود 0.6 تا 0.95 تغییر می‌کند. پراکندگی تقریباً
به نسبت \(1/\sqrt{N}\) کم می‌شود.

\textbf{تمرین 2 \ensuremath{\star} (دو شرط).} یک شرط «سرعت» (\(a = 4\)) و یک شرط «دقت»
(\(a = 10\)) با \(v = 0.2\) یکسان و هرکدام 1000 \lr{trial} شبیه‌سازی
کنید. دقت و میانگین زمان واکنش هر شرط را گزارش کنید و هر دو توزیع زمان
واکنش را رسم کنید. کدام جنبه‌ی توزیع زمان واکنش بیشتر تغییر می‌کند: کمینه،
مُد، یا دنباله؟

\textbf{تمرین 3 \ensuremath{\star}\ensuremath{\star} (سوگیری پاسخ).} قرار دهید \(a = 8\) و \(z = 5\) (شروع
نزدیک‌تر به مرز درست). \(P_c\) را تحلیلی و با شبیه‌سازی حساب کنید. بعد
میانگین زمان واکنش پاسخ‌های درست و غلط را مقایسه کنید. نتیجه را بر اساس
فاصله‌ی هر مرز از نقطه‌ی شروع توضیح دهید. \(z = 3\) چه می‌کرد؟

\textbf{تمرین 4 \ensuremath{\star}\ensuremath{\star} (بررسی یک ویژگی مدل).} برای \(z = a/2\)، با شبیه‌سازی
(20 000 \lr{trial)} نشان دهید که \emph{\lr{quantile}های} زمان واکنش
پاسخ‌های درست و غلط در حد نویز نمونه‌گیری برابرند. داده‌های واقعی اغلب در
شرایط سخت و با تأکید بر دقت، خطاهای کندتری نشان می‌دهند. آیا چنین
داده‌هایی \lr{random walk} ساده را رد می‌کنند؟ چه چیزی را تغییر می‌دادید؟

\textbf{تمرین 5 \ensuremath{\star}\ensuremath{\star}\ensuremath{\star} (پیشرفته: تغییرپذیری \lr{drift} بین \lr{trial}ها).}
\lr{\nolinkurl{simulate_trial}} را طوری تغییر دهید که در هر \lr{trial
drift} از \(v_i \sim \mathcal{N}(v, \eta)\) کشیده شود و به
\([-0.95, 0.95]\) محدود شود. از \(v = 0.2\)، \(\eta = 0.2\)، \(a = 8\)
استفاده کنید. نشان دهید که \textbf{حالا خطاها کندتر از پاسخ‌های درست‌اند}
(انتظار با 20 000 \lr{trial:} دقت \ensuremath{\approx} 0.75، حدود 12 گام برای پاسخ‌های درست
در برابر 14 گام برای خطاها) و توضیح دهید چرا. \emph{راهنمایی:} کدام
\lr{trial}ها بیشترین احتمال را دارند که به خطا ختم شوند، و آن
\lr{trial}ها چقدر طول می‌کشند؟

\subsubsection{\texorpdfstring{\lr{lab check-off} (جزو نمره‌ی \lr{lab}
حساب می‌شود)}{lab check-off (جزو نمره‌ی lab حساب می‌شود)}}

به مدرس یا دستیار آموزشی نشان دهید: (1) خروجی
\lr{\nolinkurl{check_install.py}} که همه‌ی بسته‌ها در آن
\lr{\nolinkurl{[OK]}} باشند، (2) نمودار سرعت\lr{--}دقت از بخش \lr{D}، و (3)
بهترین پارامترهای برازش و نمودار \lr{quantile}ها از بخش \lr{E.}

\begin{center}\rule{0.5\linewidth}{0.5pt}\end{center}

\subsection{\texorpdfstring{ارتباط با هوش مصنوعی و
\lr{LLM}ها}{ارتباط با هوش مصنوعی و LLMها}}

\textbf{سطوح \lr{Marr} برای یک \lr{chatbot.}} یک مدل بزرگ زبانی
\lr{(LLM)} را می‌شود در هر سه سطح \lr{Marr} توصیف کرد، و این سطوح، درست
مثل آدم‌ها، می‌توانند از هم جدا شوند. \emph{محاسباتی:} مدل آموزش دیده که
کلمه‌ی \lr{(token)} بعدی یک متن را پیش‌بینی کند. \emph{الگوریتمی:} یک
شبکه‌ی \lr{transformer} برای هر \lr{token} ممکن بعدی یک امتیاز
\lr{(«logit»)} حساب می‌کند و امتیازها را به احتمال تبدیل می‌کند.
\emph{اجرایی:} ضرب ماتریس‌ها روی \lr{GPU}ها. ادعایی مثل \lr{«LLM} سؤال را
می‌فهمد» معمولاً در سطح محاسباتی مطرح می‌شود، در حالی که شواهد (خروجی‌هایش)
را الگوریتم تولید می‌کند. وقتی چنین ادعاهایی را می‌خوانید، سطوح را از هم
جدا نگه دارید.

\textbf{یک قاعده‌ی تصمیم آشنا.} گام آخر یک \lr{LLM} یک قاعده‌ی انتخاب با
یک پارامتر نویز است، درست مثل مدل‌های شناختی این درس. احتمال \lr{token}
\(i\) یک \lr{softmax} با \textbf{دما} \lr{(temperature)} \(T\) است:
\(P(i) = e^{z_i/T}/\sum_j e^{z_j/T}\). \(T\) پایین انتخاب‌های تقریباً قطعی
و \(T\) بالا انتخاب‌های تصادفی‌تر می‌دهد. این همان قاعده‌ی \lr{softmax}
هفته‌ی 15 است، با \(\beta = 1/T\).

\begin{quote}
\textbf{\lr{Code} · کد:} \lr{\nolinkurl{code/week01.py}}, \lr{cell} 10 ·
سلول 10
\end{quote}

\textbf{خروجی مورد انتظار:}

\Needspace{8\baselineskip}
\begin{latin}
\begin{Verbatim}
T = 0.2: P(A, B, C) = [0.993 0.007 0.   ], sampled shares = [0.994 0.006 0.   ]
T = 1.0: P(A, B, C) = [0.652 0.24  0.108], sampled shares = [0.662 0.235 0.103]
T = 2.0: P(A, B, C) = [0.497 0.301 0.202], sampled shares = [0.52  0.281 0.199]
\end{Verbatim}
\end{latin}

\textbf{امتحانش کنید (بدون نیاز به \lr{API).}} - \ensuremath{\star} بزرگ‌ترین دمایی را
پیدا کنید که در آن \lr{«A»} هنوز دست‌کم 90\% مواقع انتخاب می‌شود. - \ensuremath{\star} دمای
\lr{LLM} به \emph{نویز}ی \lr{random walk} نزدیک‌تر است یا به \emph{مرز}
آن (هفته‌ی 1، بخش \lr{C)}؟ توضیح دهید. - \ensuremath{\star}\ensuremath{\star} بحث: کسی می‌گوید \lr{«chatbot}
وقتی دوباره پرسیدم جواب دیگری داد، پس جواب را نمی‌داند.» با استفاده از
این کادر توضیح دهید چه چیز دیگری ممکن است در کار باشد.

\begin{center}\rule{0.5\linewidth}{0.5pt}\end{center}

\subsection{اشتباهات رایج و برداشت‌های نادرست}

\begin{enumerate}
\def\labelenumi{\arabic{enumi}.}
\tightlist
\item
  \textbf{«مدلی که خوب برازش می‌شود درست است.»} برازش خوب لازم است ولی
  کافی نیست. یک مدل منعطف تقریباً هر چیزی، از جمله نویز، را برازش می‌کند
  (هفته‌ی 6). برازش تا وقتی ندانید مدل چه چیزی را \emph{نمی‌توانست} برازش
  کند، حرف زیادی نمی‌زند.
\item
  \textbf{«پارامتر بیشتر یعنی مدل بهتر.»} پارامتر بیشتر مدل را
  \emph{منعطف‌تر} می‌کند، که با بهتر بودن یکی نیست. انعطاف را صریحاً جریمه
  خواهیم کرد \lr{(AIC/BIC}، \lr{cross-validation}، \lr{Bayes factor).}
\item
  \textbf{قاطی کردن شبیه‌سازی با برازش.} شبیه‌سازی از پارامتر به داده
  می‌رود؛ برازش از داده به پارامتر. دانشجوها اغلب با دست‌کاری پارامترها تا
  وقتی نمودار قشنگ شود «برازش» می‌کنند. این یک \emph{کاوش} مشروع است، ولی
  روش برآورد نیست.
\item
  \textbf{پارامترها را مستقیماً قابل مشاهده دانستن.} \(v\) و \(a\)
  کمیت‌هایی \emph{استنباط‌شده‌اند} که مقدارشان به فرض‌های مدل بستگی دارد
  (مثلاً اینکه \(t_0\) معلوم است). مدل را عوض کنید، ممکن است معنای «همان»
  پارامتر هم عوض شود.
\item
  \textbf{فراموش کردن \lr{seed} تصادفی.} شبیه‌سازی بدون \lr{seed} قابل
  بازتولید نیست و هر بار که \lr{notebook} را اجرا کنید نتایج فرق می‌کند.
  از \lr{\nolinkurl{np.random.default_rng(seed)}} استفاده کنید.
\item
  \textbf{وارسی نکردن شبیه‌سازی‌ها.} باگ در شبیه‌ساز نتیجه‌گیری‌های غلطی
  می‌سازد که \emph{به نظر} معقول‌اند. با یک نتیجه‌ی تحلیلی، یک حالت حدی
  (مثلاً \(v \to 0\) باید \(P_c \to z/a\) بدهد) یا یک مقدار منتشرشده
  مقایسه کنید.
\item
  \textbf{قاطی کردن دو معنای «بیزی».} مدل‌های بیزی شناخت (ادعایی درباره‌ی
  ذهن) در برابر تحلیل داده‌ی بیزی (روشی برای پژوهشگران).
\item
  \textbf{سطوح \lr{Marr} را رقیب هم دانستن.} آن‌ها توصیف‌های مکمل یک سیستم
  واحدند، نه نظریه‌های رقیب.
\end{enumerate}

\begin{center}\rule{0.5\linewidth}{0.5pt}\end{center}

\subsection{خودآزمایی}

\textbf{1.} هرکدام را در یک جمله تعریف کنید: \emph{نظریه}، \emph{مدل
شناختی} و \emph{مدل آماری}. \textbf{پاسخ} نظریه مجموعه‌ای از ادعاهای
(معمولاً کلامی) درباره‌ی فرایندهای ذهنی است. مدل شناختی نمونه‌ی فرمال و
صریحی از چنین ادعاهایی است که می‌تواند پیش‌بینی تولید کند و پارامترهایش
کمیت‌های روان‌شناختی را نشان می‌دهند. مدل آماری ساختار داده‌ها (مثلاً میانگین
گروه‌ها، شیب‌ها) را توصیف می‌کند بی‌آنکه به یک سازوکار شناختی متعهد باشد.

\textbf{2.} هرکدام را توصیفی، فرایندی یا هنجاری دسته‌بندی کنید: (الف)
\(\mathrm{RT} = A\,N^{-b}\) به‌عنوان تابعی از تعداد \lr{trial}های تمرین
\(N\)؛ (ب) مدل \lr{random walk}؛ (ج) مشاهده‌گر ایده‌آلی که سرنخ‌های دیداری
و لامسه‌ای را با وزن قابلیت اطمینانشان ترکیب می‌کند. \textbf{پاسخ} (الف)
توصیفی: شکل داده‌ها را خلاصه می‌کند بی‌آنکه بگوید چرا. (ب) فرایندی:
سازوکاری را مشخص می‌کند که انتخاب‌ها و زمان‌های واکنش را تولید می‌کند. (ج)
هنجاری: رفتار بهینه را با توجه به اطلاعات موجود مشخص می‌کند.

\textbf{3.} مدل \lr{random walk} عمدتاً در کدام سطح \lr{Marr} قرار دارد؟
یک مثال از ادعایی درباره‌ی همان تصمیم در سطح پیاده‌سازی بزنید.
\textbf{پاسخ} سطح الگوریتمی/بازنمایی: یک رویه را توصیف می‌کند (انباشت
شواهد تا رسیدن به مرز). یک ادعای سطح پیاده‌سازی مثلاً این است: «شواهد
انباشته‌شده با نرخ شلیک نورون‌های یک ناحیه‌ی مشخص مغز حمل می‌شود و وقتی این
نرخ به یک آستانه برسد پاسخ داده می‌شود.»

\textbf{4.} تفاوت مسئله‌ی مستقیم و معکوس چیست؟ در بخش \lr{E lab} کدام را
حل کردیم؟ \textbf{پاسخ} مستقیم: با داشتن پارامترها، داده تولید (پیش‌بینی)
کنیم، یعنی شبیه‌سازی. معکوس: با داشتن داده، پارامترها را استنباط کنیم،
یعنی برازش. بخش \lr{E} یک مسئله‌ی معکوس را حل کرد (با جستجوی شبکه‌ای)، ولی
برای این کار مسئله‌ی مستقیم را بارها، یک بار برای هر نقطه‌ی شبکه، حل کرد.

\textbf{5.} با استفاده از فرمول‌ها، وقتی \(v = 0\) باشد چه بر سر \(P_c\)
می‌آید؟ (راهنمایی: به حد یا به تقارن حرکت فکر کنید.) \textbf{پاسخ} بدون
\lr{drift} حرکت متقارن است و احتمال رسیدن به مرز بالا برابر \(z/a\) است
(حد فرمول وقتی \(v \to 0\)). با شروع بی‌سوگیری، \(z = a/2\)، دقت 0.5 است،
یعنی در حد شانس. این منطقی است: بدون شواهد، تصمیم‌گیرنده حدس می‌زند.

\textbf{6.} دو گروه در دقت فرق دارند (0.80 در برابر 0.90). چرا فقط از
روی دقت نمی‌توانیم بگوییم تفاوت در احتیاط است یا در کیفیت شواهد؟ چه داده‌ی
اضافه‌ای مسئله را حل می‌کند؟ \textbf{پاسخ} هم مرز بزرگ‌تر \(a\) و هم
\lr{drift} بزرگ‌تر \(v\) دقت را بالا می‌برند، پس دقت به‌تنهایی این دو را از
هم جدا نمی‌کند. زمان واکنش مسئله را حل می‌کند: احتیاط بیشتر پاسخ‌ها را
کندتر و شواهد بهتر آن‌ها را سریع‌تر می‌کند (مثال حل‌شده، قدم 5). برازش
هم‌زمان مدل روی دقت و زمان واکنش این دو را جدا می‌کند.

\textbf{7.} چرا پیش از اعتماد به روش برازش، از پارامترهای \emph{معلوم}
داده شبیه‌سازی کردیم؟ \textbf{پاسخ} فقط با مقادیر واقعی معلوم می‌توان
بررسی کرد که آیا روش جواب درست را برمی‌گرداند (بازیابی پارامتر). اگر روشی
نتواند پارامترها را از داده‌ای که خود مدل تولید کرده بازیابی کند،
برآوردهایش از داده‌ی واقعی قابل تفسیر نیستند. این موضوع در هفته‌ی 7 کامل
بررسی می‌شود.

\begin{center}\rule{0.5\linewidth}{0.5pt}\end{center}

\subsection{تکلیف}

\subsubsection{\texorpdfstring{\lr{A0:} نصب و «ممیزی نظریه‌ی کلامی» (بدون
نمره ولی الزامی؛ موعد: ابتدای هفته‌ی
2)}{A0: نصب و «ممیزی نظریه‌ی کلامی» (بدون نمره ولی الزامی؛ موعد: ابتدای هفته‌ی 2)}}

\lr{A0} نمره ندارد ولی الزامی است: اگر \lr{A0} ناقص باشد، باید پیش از
\lr{lab check-off} هفته‌ی 2 اصلاح شود.

\textbf{بخش 1: محیط نرم‌افزاری (الزامی).} خروجی متنی
\lr{\nolinkurl{python check_install.py --sample}} را تحویل دهید که در آن
همه‌ی بسته‌ها \lr{\nolinkurl{[OK]}} باشند و میانگین پسین \lr{PyMC} دیده
شود. اگر \lr{PyMC} کار نکرد، پیام خطا و کارهایی را که امتحان کرده‌اید
تحویل دهید؛ دستیار آموزشی کمک می‌کند.

\textbf{بخش 2: ممیزی نظریه‌ی کلامی (حداکثر 1 صفحه).} یک نظریه‌ی کلامی از
درس‌های قبلی‌تان انتخاب کنید، مثلاً «سطوح پردازش»، «توجه به‌عنوان نورافکن»،
«تمرین کامل می‌کند»، «نظریه‌های دوفرایندی استدلال» یا «تداخل باعث فراموشی
می‌شود».

\begin{enumerate}
\def\labelenumi{\arabic{enumi}.}
\tightlist
\item
  نظریه را در 1 تا 2 جمله بیان کنید.
\item
  \textbf{دست‌کم پنج} تصمیمی را فهرست کنید که نظریه نامشخص گذاشته و یک
  مدل فرمال باید بگیرد (مثلاً شکل تابعی، چه چیزی پرنویز است، پارامترهای
  آزاد کدام‌اند، حالت‌های درونی چگونه به رفتار اندازه‌گیری‌شده نگاشت
  می‌شوند).
\item
  برای \textbf{دو} تا از این تصمیم‌ها یک انتخاب فرمال مشخص پیشنهاد کنید
  (یک معادله یا یک گام الگوریتمی).
\item
  \textbf{یک} پیش‌بینی را توصیف کنید که بسته به انتخاب شما فرق می‌کند، و
  آزمایشی که بتواند آن‌ها را از هم جدا کند.
\end{enumerate}

\textbf{معیارهای بازخورد (قبول / نیاز به اصلاح):}

{\def\LTcaptype{none} % do not increment counter
\begin{longtable}[]{@{}
  >{\raggedright\arraybackslash}p{(\linewidth - 4\tabcolsep) * \real{0.3333}}
  >{\raggedright\arraybackslash}p{(\linewidth - 4\tabcolsep) * \real{0.3333}}
  >{\raggedright\arraybackslash}p{(\linewidth - 4\tabcolsep) * \real{0.3333}}@{}}
\toprule\noalign{}
\begin{minipage}[b]{\linewidth}\raggedright
معیار
\end{minipage} & \begin{minipage}[b]{\linewidth}\raggedright
قبول
\end{minipage} & \begin{minipage}[b]{\linewidth}\raggedright
نیاز به اصلاح
\end{minipage} \\
\midrule\noalign{}
\endhead
\bottomrule\noalign{}
\endlastfoot
محیط نرم‌افزاری & همه‌ی بسته‌ها \lr{\nolinkurl{[OK]}}، یا گزارش روشن خطا
همراه با تلاش‌های انجام‌شده & خروجی‌ای تحویل نشده \\
نامشخص بودن نظریه & \ensuremath{\geq} 5 تصمیم متمایز و واقعاً نامشخص & کمتر از 5، یا
تکرار خود نظریه \\
فرمال‌سازی & دو انتخاب فرمال مشخص و بدون ابهام & مبهم («از یک تابع
استفاده کن») \\
پیش‌بینی متفاوت & یک پیش‌بینی مشخص به‌علاوه‌ی آزمایشی که انتخاب‌ها را از هم
جدا کند & بدون پیش‌بینی، یا پیش‌بینی‌ای که به انتخاب بستگی ندارد \\
\end{longtable}
}

بهترین ممیزی‌ها در ابتدای درس هفته‌ی 2 به‌طور مختصر بحث می‌شوند.

\begin{center}\rule{0.5\linewidth}{0.5pt}\end{center}

\subsection{مطالعه‌ی بیشتر}

\begin{itemize}
\tightlist
\item
  \textbf{\lr{Farrell} \lr{\&} \lr{Lewandowsky} (2018)، فصل‌های 1\lr{--}2.} مقدمه‌ای
  بر مدل‌سازی محاسباتی و گذر از نظریه‌های کلامی به مدل‌های فرمال.
  \emph{(منبع اصلی این هفته.)}
\item
  \textbf{\lr{Verguts} (2022)، فصل مقدمه {[}بررسی شود{]}.} انگیزه‌ای
  کوتاه و قابل فهم برای مدل‌سازی فرایندهای شناختی.
\item
  \textbf{\lr{Anderson} (2014)، فصل مقدمه {[}بررسی شود{]}.} مدل‌سازی
  شناختی را کنار علوم اعصاب محاسباتی قرار می‌دهد؛ زمینه‌ای مفید برای سطح
  پیاده‌سازی \lr{Marr.}
\item
  \emph{کلاسیک اختیاری (خارج از فهرست درس):} \lr{Marr, D.} (1982).
  \emph{\lr{Vision.}} \lr{W. H. Freeman.} منبع اصلی سه سطح تحلیل.
\end{itemize}

\emph{پیش‌نمایش هفته‌ی 2:} نظریه‌ی تشخیص علامت و منحنی‌های فراموشی را به
مدل‌های \lr{Python} تبدیل می‌کنیم و بررسی می‌کنیم هر پارامتر چگونه
پیش‌بینی‌ها را شکل می‌دهد. تکلیف \lr{A1} منتشر می‌شود.
